\documentclass[11pt]{article}

\usepackage[margin=1in]{geometry}
\usepackage{amsmath,amssymb,amsthm,mathtools}
\usepackage{booktabs}
\usepackage[hidelinks,hypertexnames=false]{hyperref}
\usepackage{microtype}

\newtheorem{theorem}{Theorem}[section]
\newtheorem{corollary}[theorem]{Corollary}
\newtheorem{proposition}[theorem]{Proposition}
\newtheorem{lemma}[theorem]{Lemma}
\theoremstyle{definition}
\newtheorem{definition}[theorem]{Definition}
\theoremstyle{remark}
\newtheorem{remark}[theorem]{Remark}

\newcommand{\N}{\mathbb N}
\newcommand{\R}{\mathbb R}
\newcommand{\C}{\mathbb C}
\newcommand{\PP}{\mathbb P}
\newcommand{\cT}{\mathcal T}
\newcommand{\Res}{\operatorname*{Res}}
\newcommand{\FP}{\operatorname*{FP}}
\newcommand{\Tr}{\operatorname{Tr}}
\newcommand{\Div}{\operatorname{Div}}
\newcommand{\norm}[1]{\lVert #1\rVert}

\title{The Determinant-Defect Completion Law\\
and Cofinal Residue-Jet Torsion}
\author{}
\date{Draft of August 4, 2026}

\begin{document}
\maketitle

\begin{abstract}
We prove a determinant-completion law independent of zeta regularization.  Let $H$ be separable
infinite-dimensional, $V$ finite-dimensional, and $p\geq2$.  Closure of graph of ordinary
finite-rank determinant pencils in Schatten and compact-open topologies is exactly
\[
\begin{aligned}
 \overline{\{(A,\det(I-A)):A:V\to\operatorname{Fin}(H)\}}
 =\bigl\{&\left(A,e^P\det{}_p(I-A)\right):\\[-2pt]
 &A:V\to\mathcal S_p,\quad
 P\in\textstyle\bigoplus_{k=1}^{p-1}\operatorname{Sym}^k(V^*)\bigr\}.
\end{aligned}
\]
Displayed parametrization is a homeomorphism, natural under pullback and direct sum, and extends
to continuous and holomorphic families.  Hence Schatten completion preserves regularized core
and forgets exactly low cyclic trace jet; every forgotten coordinate is realized by finite-rank
micro-spectra tending to zero.  We prove sharp capacity law
\[
 |\Tr(B^k)|\leq\operatorname{rank}(B)^{1-k/p}\norm{B}_p^k,
 \qquad k<p,
\]
with equality, and classify positivity constraints: positive spectra retain only linear defect,
self-adjoint spectra only linear and nonnegative quadratic defect, while graded positive spectra
recover full real jet.  For real Schatten exponent $r$, fibre jumps precisely at integer
summability walls.

Zeta determinants give geometric lifts into this completion.  If $L_n^{-1}\to A$ in
$\mathcal S_p$, Hartmann--Lesch factorization and a converse logarithmic argument show that
compact-open cluster set of normalized zeta determinants is homeomorphic to cluster set of
corrected residue jets
\[
 b_j(L_n)=-\frac1j\left(\FP_{s=j}\zeta_{L_n}(s)
 +H_{j-1}\Res_{s=j}\zeta_{L_n}(s)\right),\qquad 1\leq j<p.
\]
All limits share divisor, and divisor-cancelled relative determinant torsion is precisely
character obtained by exponentiating jet difference.  Explicit graded spectral suspensions show
all $p-1$ directions occur while Schatten limit, integer residues, and virtual divisor remain
fixed.

At threshold $p=2$ we give arithmetic realization over every number field.  For any ordering of
prime ideals, cluster set of associated finite-part constants is one closed interval in
$[e^{-\gamma},\infty]$, and every such interval occurs.  Endpoint is field-independent by exact
cancellation of Dedekind residue with Mertens constant.  Full interval is comeagre among
orderings, every interval fibre is dense, and cluster-spectrum map is nowhere continuous.
Corresponding relative zeta determinants satisfy exact prime-ideal excision and converge along a
finite-part subsequence to $e^{(1-c)z}(1-z)$; their relative torsions $e^{(d-c)z}$ realize full
real fibre although Hilbert--Schmidt limit, primary residue, divisor, and determinant line are
fixed.
\end{abstract}

\noindent\textbf{2020 Mathematics Subject Classification.}
Primary 47B10, 58J52; Secondary 30H50, 11M06, 11R42.

\noindent\textbf{Keywords.}
Determinant completion, Schatten ideal, regularized Fredholm determinant, zeta determinant,
residue jet, Wodzicki residue, determinant line, relative analytic torsion,
Euler--Kronecker constant, Mertens theorem.

\paragraph{Logical status.}
All numbered results are unconditional.  Hartmann--Lesch factorization, prime number theorem,
Mertens' theorems, continuity of regularized Fredholm determinant, and Hurwitz-zeta determinant
formula are classical inputs.  No statement assumes Riemann hypothesis.  Section~\ref{sec:boundary}
separates new cofinal theorem from established fixed-operator determinant, determinant-line,
Wodzicki-residue, and analytic-torsion theory.

\paragraph{Machine-checked core.}
Lean~4 companion supplied with source proves, without axioms or admitted goals, finite residue
algebra and cleared excision, torsion factorization and cocycle, all-order finite-difference
cancellation, nonsingularity of factorial suspension moment matrix, injective and surjective
classification of finite jet polynomials, and forward/converse compact-open convergence
criterion for factorized determinant characters after zero-free local core cancellation.  It
also proves root-of-unity micro-spectrum cancellation, exact realization of every pure
homogeneous pencil direction, and summability-transition cocycle.
Operator-specific Hartmann--Lesch factorization, Schatten determinant continuity, and Mertens
asymptotics remain cited or paper-proved analytic inputs rather than machine-checked claims.

\section{Introduction and main results}
\label{sec:intro}

Zeta determinants retain data that an underlying determinant line does not.  Quillen's metric
and Bismut--Freed connection make this distinction geometric; their local curvature defects are
governed by residue terms, while relative determinants and analytic torsion retain global finite
parts and gluing data~\cite{Quillen1985,BismutFreed1986,PaychaRosenberg2000,Scott1999}.
At the level of pseudodifferential operators, Wodzicki residue produces a local multiplicative
determinant, whereas zeta determinant is nonlocal and generally nonmultiplicative
\cite{KontsevichVishik1994,Scott2005}.  Two types of information cannot be interchanged.

Quillen anticipated a principal bundle of determinant trivializations under exponentials of
bounded-degree polynomial functions.  In a local elliptic setting, Dang proved this picture:
renormalized determinants obtained from local counterterms form a torsor under local polynomial
functionals and give holomorphic trivializations of determinant line~\cite{Dang2022}.  Our
question is orthogonal to locality and precedes choice of renormalization: which polynomial
ambiguities are forced merely by closing \emph{uncorrected finite-dimensional determinants} in
Schatten topology?  The answer is exact.  Every low cyclic trace discarded by that topology
becomes an independent fibre coordinate, and there are no others.  This topological completion
law supplies common source from which corrected zeta jets, relative torsion characters, and
graded anomaly coordinates below are obtained by imposing geometric or arithmetic constraints.

Parent law is already visible before zeta regularization.  Let $H$ be separable
infinite-dimensional Hilbert space, let $V$ be finite-dimensional complex vector space, and
write $\operatorname{Fin}(H)$ for finite-rank operators.  For linear pencil
$A:V\to\operatorname{Fin}(H)$ put
\[
 D_A(v)=\det(I-A(v)),\qquad D_{p,A}(v)=\det{}_p(I-A(v)).
\]
Let $\operatorname{Hol}_1(V)$ carry compact-open topology and consist of entire functions taking
value one at origin.  Give $\operatorname{Hom}(V,\mathcal S_p)$ any operator norm induced by
norm on $V$; all choices are equivalent.  Finally put
\[
 J_{<p}(V)=\bigoplus_{k=1}^{p-1}\operatorname{Sym}^k(V^*).
\]

\begin{theorem}[Determinant-defect completion law]
\label{thm:completion-law}
For every integer $p\geq2$, closure of finite-rank determinant graph in
$\operatorname{Hom}(V,\mathcal S_p)\times\operatorname{Hol}_1(V)$ is
\begin{equation}
 \overline{\left\{(A,D_A):A\in\operatorname{Hom}(V,\operatorname{Fin}(H))\right\}}
 =
 \left\{\left(A,e^P D_{p,A}\right):
 A\in\operatorname{Hom}(V,\mathcal S_p),\ P\in J_{<p}(V)\right\}.       \tag{DC}
\end{equation}
Map
\begin{equation}
\Theta_{p,V}:\operatorname{Hom}(V,\mathcal S_p)\times J_{<p}(V)
 \longrightarrow\overline{\operatorname{Graph}(D)},\qquad
 (A,P)\longmapsto(A,e^P D_{p,A})                                  \tag{1.D}
\end{equation}
is homeomorphism.  It intertwines linear pullback $W\to V$, unitary conjugacy on $H$, and
orthogonal direct sum; under direct sum, jet coordinates add.  Hence every fibre is principal
homogeneous space for natural group $\exp J_{<p}(V)$.
\end{theorem}

\begin{theorem}[Parametrized and holomorphic completion]
\label{thm:parametrized-completion}
Let $X$ be compact Hausdorff, and let $C_{\rm fs}$ denote continuous pencils whose operators
have one common finite-dimensional support.  Take closure in
$C(X,\operatorname{Hom}(V,\mathcal S_p))\times C(X,\operatorname{Hol}_1(V))$, with uniform
$\mathcal S_p$ topology in $x$ and compact-open topology in $(x,v)$.  Then
\begin{align}
&\overline{\left\{\left(A,(x,v)\mapsto\det(I-A_x(v))\right):
 A\in C_{\rm fs}\!\left(X,\operatorname{Hom}(V,\operatorname{Fin}(H))\right)\right\}} \notag\\
&\quad=\left\{\left(A,(x,v)\mapsto e^{P_x(v)}D_{p,A_x}(v)\right):
 A\in C\!\left(X,\operatorname{Hom}(V,\mathcal S_p)\right),
 P\in C\!\left(X,J_{<p}(V)\right)\right\}.                 \tag{1.P}
\end{align}
Parametrization by $(A,P)$ is homeomorphism from
\begin{equation*}
 C\!\left(X,\operatorname{Hom}(V,\mathcal S_p)\right)
 \times C\!\left(X,J_{<p}(V)\right).
\end{equation*}
Same statement holds sheafwise for holomorphic
families over any complex manifold, with compact-open topology on base.  Construction commutes
with restriction and holomorphic pullback.  Thus completion fibre over a family is gauge group
of continuous, respectively holomorphic, sections of $\exp J_{<p}(V)$.
\end{theorem}

These finite-support families are dense in continuous $\mathcal S_p$ families over compact
bases.

This is not Hadamard factorization: right side is \emph{closure forced by Schatten completion},
and every fibre direction is realized by determinants of finite-rank pencils converging to same
operator pencil.  Nor is it fixed-operator zeta--Fredholm factorization: statement contains no
zeta function.  It identifies complete nonclosability relation of ordinary determinant.

\begin{theorem}[Sharp micro-spectrum capacity and superization]
\label{thm:micro-spectrum}
Let $1\leq k<p$ and let $B$ have finite rank $R$.  Then
\begin{equation}
 |\Tr(B^k)|\leq R^{1-k/p}\norm{B}_p^k.                    \tag{1.E}
\end{equation}
For every $c\in\C$ equality is attained by a normal root-of-unity block with
$\Tr(B^q)=0$ for $1\leq q<k$ and $\Tr(B^k)=c$.  Consequently a degree-$k$ defect $c\neq0$
carried with Schatten energy $\norm{B}_p=\varepsilon$ requires and permits sharp rank scale
\begin{equation}
 R=|c|^{p/(p-k)}\varepsilon^{-pk/(p-k)}
 \quad\text{up to integrality of }R.                      \tag{1.R}
\end{equation}
No nonzero degree-$p$ moment survives a perturbation tending to zero in $\mathcal S_p$.

At zero fibre, possible convergent low-moment defects have following exact statistics:
\begin{center}
\begin{tabular}{ll}
\toprule
allowed finite-rank micro-spectrum & limiting moments $(m_1,\ldots,m_{p-1})$\\
\midrule
complex & all of $\C^{p-1}$\\
positive & $(m_1,0,\ldots,0)$, $m_1\geq0$\\
self-adjoint & $(m_1,m_2,0,\ldots,0)$, $m_1\in\R$, $m_2\geq0$\\
graded positive, using supertrace & all of $\R^{p-1}$\\
\bottomrule
\end{tabular}
\end{center}
Rows are truncated when $p=2$ or $3$.  Thus positivity alone loses higher hierarchy, whereas
standard $\mathbb Z/2$-graded Grothendieck completion recovers every real defect direction.
\end{theorem}

\begin{corollary}[Summability phase and anomaly filtration]
\label{cor:summability-phase}
Let $1\leq r<\infty$ be real and $m=\lceil r\rceil$.  Replace $\mathcal S_p$ by
$\mathcal S_r$, $\det_p$ by $\det_m$, and $J_{<p}(V)$ by
\[
 J_{<r}(V):=\bigoplus_{\substack{k\in\N\\1\leq k<r}}
 \operatorname{Sym}^k(V^*).
\]
Then (DC) and homeomorphism (1.D) remain valid.  Degree-$k$ capacity exponent is
$r/(r-k)$.  Hence defect fibre changes only when $r$ crosses an integer; passage from
$\mathcal S_p$ to any $\mathcal S_r$ with $p<r\leq p+1$ adds exactly
$\operatorname{Sym}^p(V^*)$.

For integers $2\leq p<q$, completions form natural tower.  On
$A\in\operatorname{Hom}(V,\mathcal S_p)$ put
\[
 T_{p,q}(A)(v)=\sum_{j=p}^{q-1}\frac1j\Tr(A(v)^j).
\]
Embedding preserving completed determinant character is
\begin{equation}
 \iota_{p,q}(A,P)=\bigl(A,P-T_{p,q}(A)\bigr),\qquad
 e^P D_{p,A}=e^{P-T_{p,q}(A)}D_{q,A}.                    \tag{1.S}
\end{equation}
These embeddings satisfy $\iota_{q,r}\iota_{p,q}=\iota_{p,r}$; associated new anomaly layer at
$p$ is $\operatorname{Sym}^p(V^*)$.
\end{corollary}

Theorems~\ref{thm:completion-law} and \ref{thm:micro-spectrum} are proved in
Section~\ref{sec:completion}.  They explain every result below.  Zeta regularization selects
canonical lift $(L^{-1},P_L)$ in completion relation; residue correction computes its jet
coordinate; relative torsion is transition character in principal fibre; graded suspension is
positive realization of full fibre; prime-ideal orderings describe constrained arithmetic
dynamics inside one real line of fibre.

\begin{center}
\small
\begin{tabular}{p{0.28\textwidth}p{0.31\textwidth}p{0.28\textwidth}}
\toprule
completion datum & zeta/geometric shadow & arithmetic collateral\\
\midrule
$A\in\mathcal S_p$ & inverse $L^{-1}$ & prime-sieved inverse $Q_{K,P}$\\
$\exp J_{<p}(V)$ & corrected residue jet $b(L)$ & finite parts $\kappa_K(P)$\\
principal-fibre transition & relative determinant/torsion & $e^{(d-c)z}$\\
$\operatorname{Sym}^k(V^*)$ layer & order-$k$ residue correction & colored suspension coordinate\\
graded positive group completion & graded determinant/analytic torsion & signed scale suspension\\
compact-open graph closure & cofinal determinant clusters & interval spectrum of orderings\\
\bottomrule
\end{tabular}
\end{center}

Write $H_0=0$ and $H_j=1+\cdots+1/j$.  For integer $p\geq2$, call positive invertible
$L$ an admissible $p$-datum if $L^{-1}\in\mathcal S_p$, its spectral zeta function is regular
at zero and meromorphic with at most simple poles at $1,\ldots,p-1$, and standard
discrete-dimension-spectrum hypotheses of Hartmann--Lesch hold.  Put
\begin{equation}
 b_j(L)=-\frac1j\left(\FP_{s=j}\zeta_L(s)
     +H_{j-1}\Res_{s=j}\zeta_L(s)\right),\qquad
 P_{b(L)}(z)=\sum_{j=1}^{p-1}b_j(L)z^j.                 \tag{1.0}
\end{equation}
At regular $j$, finite part means $\zeta_L(j)$.  For classical elliptic operators residue in
(1.0) is local Wodzicki residue; finite part is global.

\begin{theorem}[Universal cofinal residue-jet law]
\label{thm:universal-law}
Let $L_n$ be admissible $p$-data and suppose $A_n=L_n^{-1}\to A$ in $\mathcal S_p$.  Set
\[
 \mathcal D_n(z)=\frac{\det_\zeta(L_n-z)}{\det_\zeta L_n},\qquad
 \Phi_A(b)(z)=e^{P_b(z)}\det{}_p(I-zA).
\]
Then $\Phi_A:\R^{p-1}\to\operatorname{Hol}_1(\C)$ is topological embedding for compact-open
topology, where subscript means value one at zero, and
\[
 \operatorname{Clust}_{\rm co}(\mathcal D_n)
 =\Phi_A\!\left(\operatorname{Clust}_{\R^{p-1}}(b(L_n))\right). \tag{1.U}
\]
In particular, locally uniform convergence of $\mathcal D_{n_k}$ occurs if and only if
$b(L_{n_k})$ converges, and every determinant limit has divisor of $\det_p(I-zA)$.  For two
finite jet limits $b,c$, divisor-cancelled relative torsion is
\[
 \mathcal T_{b,c}(z)=\frac{\Phi_A(b)(z)}{\Phi_A(c)(z)}
 =e^{P_b(z)-P_c(z)},\qquad
 \mathcal T_{b,c}\mathcal T_{c,d}=\mathcal T_{b,d}.      \tag{1.T}
\]
If jet cluster set has nonempty interior in $\R^{p-1}$, generated torsion group is full group of
characters $\exp(\sum_{j=1}^{p-1}v_jz^j)$, $v\in\R^{p-1}$.
\end{theorem}

\begin{theorem}[Full-jet sharpness by spectral suspension]
\label{thm:full-jet}
For every $p\geq2$ there is explicit countable graded family of positive diagonal admissible
$p$-data, indexed by colored primes, with following properties along every colored-prime
exhaustion: every constituent inverse converges in $\mathcal S_p$ to fixed rank-one operator,
integer residue vector $(\Res_{s=k}\zeta(s))_{1\leq k<p}$ tends to zero, and virtual determinant
divisor converges.  Nevertheless, on comeagre set of colored-prime orderings, corrected virtual
jet cluster set is full product of $p-1$ nondegenerate intervals and has nonempty interior.
Consequently group generated by divisor-cancelled relative determinant torsions is
\[
 \left\{\exp\!\left(\sum_{k=1}^{p-1}v_kz^k\right):v\in\R^{p-1}\right\}. \tag{1.F}
\]
All $p-1$ jet coordinates are necessary: deleting any one makes corresponding character
direction unrecoverable.
\end{theorem}

Fixed-operator polynomial factorization behind theorem is established
in~\cite{HartmannLesch2022}.  New assertion is cofinal completeness: no discontinuity beyond
finite-part residue jet is possible, none of jet can be discarded, and cluster topology and
relative torsion are transported exactly.  Number-field theorem below gives natural sharp
realization with fixed limiting operator, residue, and divisor but maximal jet ambiguity.

Let $p_1<p_2<\cdots$ denote natural ordering of primes and put
\[
 E_r=\{p_1,\ldots,p_r\},\qquad E_0=\varnothing.
\]
For finite $P\subset\PP$, define
\begin{equation}
 a(P)=\prod_{p\in P}\left(1-\frac1p\right),\qquad
 b(P)=\gamma+\sum_{p\in P}\frac{\log p}{p-1},\qquad
 \gamma(P)=a(P)b(P).                                      \tag{1.1}
\end{equation}
This is Diamond--Ford generalized Euler constant~\cite{DiamondFord2008}.  Rational case will
serve as a transparent proof model before Section~\ref{sec:numberfields} establishes number-field
theorem.  Equivalently,
\begin{equation}
 \gamma(P)=\lim_{x\to\infty}
 \left(
 \sum_{\substack{m\leq x\\(m,M_P)=1}}\frac1m-a(P)\log x
 \right),\qquad M_P=\prod_{p\in P}p.                     \tag{1.2}
\end{equation}

Let $\sigma=(q_1,q_2,\ldots)$ be a permutation of $\PP$ and
$P_n^\sigma=\{q_1,\ldots,q_n\}$.  Cluster sets are taken in order compactification
$[0,\infty]$.

\begin{theorem}[Cofinal torsion-spectrum classification]
\label{thm:classification}
For every permutation $\sigma$ of primes, cluster set
\[
 C(\sigma)=\operatorname{Clust}_{n\to\infty}\gamma(P_n^\sigma)
\]
is a nonempty closed interval in $[e^{-\gamma},\infty]$.  Conversely, for every
\[
 e^{-\gamma}\leq\alpha\leq\beta\leq\infty
\]
there is a permutation $\sigma$ with
\begin{equation}
 C(\sigma)=[\alpha,\beta].                                \tag{1.3}
\end{equation}
Degenerate interval $[\infty,\infty]=\{\infty\}$ is allowed.
\end{theorem}

Lower endpoint is cofinal, not unrestricted finite-set infimum.  Diamond--Ford proved that all
finite-set values are dense in $[\Gamma,\infty)$, where
$\Gamma=\inf_P\gamma(P)$, and that $\Gamma<e^{-\gamma}$ if RH is false.  Theorem
\ref{thm:classification} shows values below $e^{-\gamma}$, if they exist, cannot persist along
any nested exhaustion of primes.

On
\[
 H_P=\ell^2\{m\geq1:(m,M_P)=1\}
\]
let $D_Pe_m=me_m$ and $Q_P=D_P^{-1}$.  Its spectral zeta function is
\begin{equation}
 Z_P(s)=\Tr(Q_P^s)=\zeta(s)\prod_{p\in P}(1-p^{-s}).       \tag{1.4}
\end{equation}

\begin{theorem}[Residue-corrected relative determinant]
\label{thm:determinant}
For every finite $P$,
\begin{equation}
 \Res_{s=1}Z_P(s)=a(P),\qquad
 \FP_{s=1}Z_P(s)=\gamma(P).                               \tag{1.5}
\end{equation}
Locally uniform limit
\begin{equation}
 \Delta_P(z)=\lim_{N\to\infty}N^{a(P)z}
 \prod_{\substack{1\leq m\leq N\\(m,M_P)=1}}
 \left(1-\frac zm\right)                                 \tag{1.6}
\end{equation}
exists on $\C$ and satisfies
\begin{align}
 \Delta_P(z)
 &=e^{-\gamma(P)z}\det{}_2(I-zQ_P),                       \tag{1.7}\\
 &=M_P^{a(P)z}
 \prod_{\substack{1\leq r\leq M_P\\(r,M_P)=1}}
 \frac{\Gamma(r/M_P)}{\Gamma((r-z)/M_P)},                \tag{1.8}\\
 &=\frac{\det_\zeta(D_P-z)}{\det_\zeta D_P}.             \tag{1.9}
\end{align}
If $p\notin P$, then
\begin{equation}
 \boxed{
 \Delta_{P\cup\{p\}}(z)
 =p^{-a(P)z/p}\frac{\Delta_P(z)}{\Delta_P(z/p)}.}         \tag{1.10}
\end{equation}
The quotient in (1.10) has its removable singularities filled.  Equivalently, the global
entire identity is
\[
 p^{a(P)z/p}\Delta_{P\cup\{p\}}(z)\Delta_P(z/p)=\Delta_P(z).
\]
\end{theorem}

\begin{theorem}[Determinant clusters and relative torsion]
\label{thm:clusters}
Let $\sigma$ be a permutation of primes.  If
$\gamma(P_{n_j}^\sigma)\to c<\infty$, then, locally uniformly on $\C$,
\begin{equation}
 \Delta_{P_{n_j}^\sigma}(z)\longrightarrow
 \Delta_c(z):=e^{(1-c)z}(1-z).                            \tag{1.11}
\end{equation}
Conversely, every locally uniformly convergent subsequence of
$\Delta_{P_n^\sigma}$ has form \emph{(1.11)} for a finite cluster value $c$.
Thus finite determinant cluster set is exactly
\begin{equation}
 \{e^{(1-c)z}(1-z):c\in C(\sigma)\cap\R\}.                \tag{1.12}
\end{equation}
For finite $c,d$ relative determinant torsion is
\begin{equation}
 \cT_{c,d}(z)=\frac{\Delta_c(z)}{\Delta_d(z)}
 =e^{(d-c)z},\qquad
 \cT_{c,d}\cT_{d,e}=\cT_{c,e}.                           \tag{1.13}
\end{equation}
Here and below quotient of determinant sections means unique holomorphic extension after
cancelling their common divisor; equivalently
$\Delta_c=\cT_{c,d}\Delta_d$ everywhere.  This convention removes apparent $0/0$ at $z=1$.
All limiting determinants have same divisor $[1]$, while pairwise torsions realize every real
logarithmic slope.
\end{theorem}

\begin{theorem}[Number-field universality]
\label{thm:numberfield-main}
Let $K$ be any number field and order its nonzero prime ideals arbitrarily.  For initial sets
$P_n$ in this ordering, define
\[
 \kappa_K(P)=\FP_{s=1}\left[
 \zeta_K(s)\prod_{\mathfrak p\in P}(1-N\mathfrak p^{-s})\right].
\]
Cluster sets of $\kappa_K(P_n)$ are exactly nonempty closed intervals
$[\alpha,\beta]\subset[e^{-\gamma},\infty]$.  For the ideal-norm operator $D_{K,P}$, all finite
determinant cluster germs are $e^{(1-c)z}(1-z)$ with $c$ in corresponding interval, and all
relative torsions are $e^{(d-c)z}$.  Hilbert--Schmidt limit, primary residue, and divisor are
independent of ordering.
\end{theorem}

\begin{theorem}[Generic maximal torsion and nowhere continuity]
\label{thm:generic}
Give space $\mathfrak S_K$ of all orderings of $\mathcal P_K$ pointwise-convergence topology.
For every nonempty closed interval $I\subset[e^{-\gamma},\infty]$, fibre
\[
 \{\sigma\in\mathfrak S_K:C_K(\sigma)=I\}
\]
is dense.  Full-spectrum fibre $C_K(\sigma)=[e^{-\gamma},\infty]$ is comeagre; every proper
interval fibre is dense and meagre.  Map $\sigma\mapsto C_K(\sigma)$ to hyperspace of nonempty
compact subsets is nowhere continuous.  On comeagre full-spectrum fibre, determinant cluster
germs contain every $e^{(1-c)z}(1-z)$ with finite $c\geq e^{-\gamma}$ and relative-torsion slopes
are all of $\R$.
\end{theorem}

\paragraph{Organization.}
Sections~\ref{sec:finite}--\ref{sec:generic} first develop constrained arithmetic and zeta
models, making visible defects that survive fixed operator, residue, and divisor.  Section
\ref{sec:geometry} identifies determinant-line and torsion meaning.  Section
\ref{sec:completion} then proves topology-forced parent law, its family form, sharp capacity,
positivity classification, and summability filtration.  This concrete-to-structural order also
keeps classical fixed-operator inputs separate from new graph-closure theorem.  Final sections
state exact novelty boundary and consequences.

\section{Finite-set identities and extremality}
\label{sec:finite}

We first isolate exact one-prime dynamics.

\begin{lemma}[One-prime update]
\label{lem:update}
If $p\notin P$, then
\begin{equation}
 \gamma(P\cup\{p\})-\gamma(P)
 =\frac{a(P)}p\bigl(\log p-b(P)\bigr).                    \tag{2.1}
\end{equation}
In particular, adding $p$ raises $\gamma(P)$ exactly when $\log p>b(P)$.
\end{lemma}

\begin{proof}
From (1.1),
\[
 a(P\cup\{p\})=a(P)\left(1-\frac1p\right),\qquad
 b(P\cup\{p\})=b(P)+\frac{\log p}{p-1}.
\]
Multiplication and subtraction give (2.1).
\end{proof}

\begin{lemma}[Relative Diamond--Ford extremality]
\label{lem:relative-extremal}
If $E_k\subset P$, then
\begin{equation}
 \gamma(P)\geq\inf_{r\geq k}\gamma(E_r).                 \tag{2.2}
\end{equation}
\end{lemma}

\begin{proof}
We repeat descent underlying Diamond--Ford extremal theorem and retain relative constraint.
Suppose finite set $Q$ is not an initial prime segment and let $t$ be its largest prime.  Put
$Q'=Q\setminus\{t\}$.  If $\gamma(Q')<\gamma(Q)$, replace $Q$ by $Q'$.  Otherwise let $u$ be
smallest prime not in $Q$ and replace $Q$ by $Q'\cup\{u\}$.

For prime $v\notin Q'$,
\[
 \gamma(Q'\cup\{v\})
 =\gamma(Q')f(v),\qquad
 f(v)=1-\frac1v+\frac{\log v}{vA},\qquad A=b(Q').
\]
Function $f$ increases on $(1,e^{A+1})$, decreases on $(e^{A+1},\infty)$, and tends to $1$ at
infinity.  In second case $f(t)\leq1$, hence $u<t\leq e^{A+1}$ and $f(u)<f(t)$.  Therefore both
allowed moves strictly decrease $\gamma$.

Process terminates at an initial segment $E_r$.  If $E_k\subset Q$ at start, every nonterminal
$Q$ has largest element $t>p_k$; removing $t$ preserves $E_k$, and smallest missing prime $u$ is
also larger than $p_k$.  Relative constraint is preserved and terminal index satisfies $r\geq k$.
Thus $\gamma(P)\geq\gamma(E_r)$ for some $r\geq k$, proving (2.2).
\end{proof}

Mertens' theorems and prime number theorem give
\begin{equation}
 a(E_r)\sim\frac{e^{-\gamma}}{\log p_r},\qquad
 b(E_r)=\log p_r+O(1),\qquad
 \gamma(E_r)\longrightarrow c_*:=e^{-\gamma}.             \tag{2.3}
\end{equation}

\begin{corollary}[Cofinal lower bound]
\label{cor:lower}
For every increasing exhaustion $P_1\subset P_2\subset\cdots$ of primes,
\begin{equation}
 \liminf_{n\to\infty}\gamma(P_n)\geq c_*.                \tag{2.4}
\end{equation}
\end{corollary}

\begin{proof}
For fixed $k$, eventually $E_k\subset P_n$.  Lemma~\ref{lem:relative-extremal} gives
$\gamma(P_n)\geq\inf_{r\geq k}\gamma(E_r)$.  Let first $n$ and then $k$ tend to infinity, using
(2.3).
\end{proof}

This explains why unrestricted lower endpoint $\Gamma$ disappears under cofinality.  It also
shows result is unconditional: possible RH-dependent dips of $\gamma(E_r)$ below $c_*$ vanish
in every tail infimum.

\section{Local mesh and prime-block controls}
\label{sec:controls}

\begin{lemma}[Vanishing mesh on bounded windows]
\label{lem:mesh}
Let $P_j$ be finite prime sets and let $q_j\notin P_j$ be primes with $q_j\to\infty$.  If
$\gamma(P_j)\leq M$, then
\begin{equation}
 \left|\gamma(P_j\cup\{q_j\})-\gamma(P_j)\right|
 \leq\frac{M+\log q_j}{q_j}.                              \tag{3.1}
\end{equation}
Consequently one-prime jumps tend to zero whenever initial values stay in a fixed bounded window.
\end{lemma}

\begin{proof}
If right side of (2.1) is negative, its absolute value is at most
$a(P_j)b(P_j)/q_j=\gamma(P_j)/q_j$.  If positive, it is at most
$a(P_j)\log q_j/q_j\leq\log q_j/q_j$.  Adding bounds proves (3.1).
\end{proof}

For $u\geq0$ and $y>1$, write
\begin{equation}
 B(y,u)=\{p\in\PP:y<p\leq e^u y\}.                        \tag{3.2}
\end{equation}

\begin{lemma}[Localized upward control]
\label{lem:up}
Fix finite $P$ and $u\geq0$.  As $y\to\infty$,
\begin{align}
 \prod_{p\in B(y,u)}\left(1-\frac1p\right)&\longrightarrow1,
                                                                    \tag{3.3}\\
 \sum_{p\in B(y,u)}\frac{\log p}{p-1}&\longrightarrow u,           \tag{3.4}\\
 \gamma(P\cup B(y,u))&\longrightarrow\gamma(P)+a(P)u.              \tag{3.5}
\end{align}
Moreover, $y$ may be chosen so that, when primes in $B(y,u)$ are inserted in increasing order,
every step raises $\gamma$ and largest step is arbitrarily small.
\end{lemma}

\begin{proof}
Mertens' estimate for prime reciprocals gives
\[
 \sum_{y<p\leq e^uy}\frac1p
 =\log\log(e^uy)-\log\log y+o(1)=o(1).
\]
Together with $\sum_{p>y}p^{-2}=o(1)$, this proves (3.3).  Partial summation from prime number
theorem gives the Mertens asymptotic
\[
 \sum_{p\leq x}\frac{\log p}{p}=\log x+C+o(1)
\]
for an absolute constant $C$.  Since
\[
 \sum_{p>y}\left|\frac{\log p}{p-1}-\frac{\log p}{p}\right|=o(1),
\]
(3.4) follows.  Equations (3.3), (3.4), and (1.1) imply (3.5).

During block, $b$-coordinate is at most $b(P)+u+1$ for large $y$.  Choose
$\log y>b(P)+u+1$.  Lemma~\ref{lem:update} makes every step positive.  Its size is at most
$\log(e^uy)/y$, which tends to zero.
\end{proof}

\begin{lemma}[Downward cofinal control]
\label{lem:down}
Let $P$ be finite.  Inserting missing primes in natural order eventually produces an initial
segment $E_r$, and such terminal segments have $\gamma(E_r)\to c_*$.  More precisely, fix
$\alpha>c_*$.  Once least missing prime is sufficiently large, every least-missing-prime
insertion made while $\gamma(P)\geq\alpha$ strictly decreases $\gamma(P)$.
\end{lemma}

\begin{proof}
If $X\geq\max P$, adjoining all primes at most $X$ yields exactly initial segment
$\{p:p\leq X\}$, so first assertion follows from (2.3).

Suppose $q=p_k$ is least missing prime.  Then $E_{k-1}\subset P$, hence
$a(P)\leq a(E_{k-1})$.  By (2.3),
\[
 a(E_{k-1})\log p_k\longrightarrow c_*.
\]
For all large $k$, $a(P)\log q<\alpha$.  If $\gamma(P)\geq\alpha$, then
$b(P)=\gamma(P)/a(P)>\log q$, and Lemma~\ref{lem:update} gives strict decrease.
\end{proof}

For endpoint $\alpha=c_*$, same proof shows least-missing insertions decrease whenever
$\gamma(P)\geq c_*+\varepsilon$, once least missing prime is large in terms of $\varepsilon$.
Thus downward sweep may oscillate inside a shrinking neighborhood of $c_*$, but cannot make a
fixed positive excursion above it.

\section{Necessity: every cofinal cluster set is an interval}
\label{sec:necessity}

Let $x_n=\gamma(P_n^\sigma)$.  Since $\sigma$ is a permutation, $q_n\to\infty$: below every fixed
bound there are only finitely many primes, and all occur at finite indices.
Corollary~\ref{cor:lower} gives
\begin{equation}
 \liminf x_n\geq c_*.                                     \tag{4.1}
\end{equation}

Cluster set is nonempty and closed in compact space $[0,\infty]$.  It remains to prove order
convexity.  Let $u<v$ be two finite cluster values and choose infinitely many passages of sequence
from a neighborhood of $u$ to a neighborhood of $v$, or conversely.  Fix $t\in(u,v)$.  At first
crossing of $t$, both endpoints are eventually bounded: for upward crossing this is immediate,
while for downward crossing
\begin{equation}
 x_{n+1}=x_n\left(1-\frac1{q_{n+1}}\right)
 +\frac{a(P_n^\sigma)\log q_{n+1}}{q_{n+1}}
 \geq x_n\left(1-\frac1{q_{n+1}}\right),                  \tag{4.2}
\end{equation}
so bounded terminal endpoint forces bounded initial endpoint.  Lemma~\ref{lem:mesh} makes crossing
mesh tend to zero.  Some crossing terms converge to $t$, and $t$ is a cluster value.

Same crossing argument applies when $v=\infty$ and $u<\infty$: every finite $t>u$ is then a
cluster value.  If there is no finite cluster value, compactness forces cluster set to be
$\{\infty\}$.  Combined with (4.1), this proves necessity half of
Theorem~\ref{thm:classification}.

\section{Sufficiency: realization of every interval}
\label{sec:realization}

We construct an ordering as concatenation of finite blocks of previously unused primes.  At every
downward stage we insert least missing primes; therefore least missing prime tends to infinity and
resulting infinite list exhausts $\PP$.  Upward blocks are always chosen above all primes used so
far, so blocks are disjoint.

First let
\[
 c_*\leq\alpha<\beta<\infty
\]
and choose $\varepsilon_j\downarrow0$.  A natural initial segment starts arbitrarily close to
$c_*$.  If $\alpha>c_*$, one application of Lemma~\ref{lem:up} raises it to an arbitrarily small
neighborhood of $\alpha$; if $\alpha=c_*$, no raising is needed.  After this finite
initialization, perform following cycle.

\begin{enumerate}
\item Starting within $\varepsilon_j$ of $\alpha$, apply Lemma~\ref{lem:up} with target
$\beta-\varepsilon_j/2$.  Choose block far enough that terminal value is within
$\varepsilon_j/2$ of target and every step is smaller than $\varepsilon_j$.  This is an
increasing sweep from lower to upper endpoint.
\item Insert least missing primes until value first enters
$\varepsilon_{j+1}$-neighborhood of $\alpha$.  If $\alpha>c_*$,
Lemma~\ref{lem:down} makes late part of sweep decreasing until first crossing.  If
$\alpha=c_*$, fill through a sufficiently large natural cutoff; (2.3) reaches desired
neighborhood, and endpoint version of Lemma~\ref{lem:down} prevents a fixed excursion.
Lemma~\ref{lem:mesh} makes overshoot smaller than $\varepsilon_{j+1}$.
\end{enumerate}

Every prime is eventually used.  All sufficiently late terms lie in an arbitrarily small
neighborhood of $[\alpha,\beta]$.  Conversely, each upward sweep has mesh tending to zero and
crosses every point of $[\alpha,\beta]$.  Cluster set is exactly that interval.

For singleton $[c,c]$ with $c>c_*$, consume at least one least missing prime and continue until
value lies just below $c$, then use a far block to raise it to within $\varepsilon_j$ of $c$.
Lemmas~\ref{lem:down}
and~\ref{lem:up} show every term in $j$th cycle lies within $O(\varepsilon_j)$ of $c$.  Natural
ordering realizes $[c_*,c_*]$.

For $[\alpha,\infty]$ with $\alpha<\infty$, use same construction with upper targets
$\beta_j\to\infty$.  Upward sweeps have vanishing mesh on every fixed bounded interval, so every
finite $t\geq\alpha$ is a cluster value; downward control excludes finite cluster values below
$\alpha$, and upper targets supply $\infty$.

Finally, to realize $\{\infty\}$, suppose least missing prime at stage $j$ is $q_j$.  Begin stage
with a value at least $2j$.  Insert $q_j$.  By (4.2), new value is at least
$2j(1-1/q_j)$, hence at least $j$ for all sufficiently large $j$.  A far upward block then raises
value above $2(j+1)$.  All later terms tend to infinity, while one least missing prime is consumed
at every stage.  This completes proof of Theorem~\ref{thm:classification}.

\section{Universal residue-jet law and residue-corrected determinants}
\label{sec:determinants}

We first prove Theorem~\ref{thm:universal-law}.  For $|z|$ below spectrum, relative spectral
zeta function has binomial expansion
\begin{equation}
 \Tr\bigl((L-z)^{-s}-L^{-s}\bigr)
 =\sum_{j\geq1}\frac{(s)_j}{j!}z^j\zeta_L(s+j).          \tag{6.0a}
\end{equation}
At a simple pole $s=j$ write
$\zeta_L(j+s)=r_j/s+\kappa_j+O(s)$.  Since
\[
 \frac{(s)_j}{j!}=\frac{s}{j}+\frac{H_{j-1}}j s^2+O(s^3),
\]
$j$th coefficient in negative derivative at zero of (6.0a) is
$-(\kappa_j+H_{j-1}r_j)/j=b_j(L)$.  Same formula holds at regular point with $r_j=0$.
Terms $j\geq p$ sum to logarithm of $p$-regularized Fredholm determinant.  Thus established
zeta--Fredholm formula~\cite{HartmannLesch2022} becomes, with present sign convention,
\begin{equation}
 \frac{\det_\zeta(L-z)}{\det_\zeta L}
 =e^{P_{b(L)}(z)}\det{}_p(I-zL^{-1}).                    \tag{6.0b}
\end{equation}

Now let $A_n=L_n^{-1}\to A$ in $\mathcal S_p$.  Continuity estimate for regularized Fredholm
determinants~\cite{Simon2005} gives
\begin{equation}
 F_n(z):=\det{}_p(I-zA_n)\longrightarrow
 F(z):=\det{}_p(I-zA)                                   \tag{6.0c}
\end{equation}
locally uniformly.  Therefore $b(L_{n_k})\to b$ implies
$\mathcal D_{n_k}\to\Phi_A(b)$ locally uniformly.

Conversely suppose $\mathcal D_{n_k}\to G$ locally uniformly.  All functions in sight equal one
at zero.  Hence on one fixed disc about zero, $G,F$, and, eventually, $\mathcal D_{n_k},F_{n_k}$
are zero-free.  Choose logarithms vanishing at zero.  Equations (6.0b)--(6.0c) give
\begin{equation}
 P_{b(L_{n_k})}=\log\mathcal D_{n_k}-\log F_{n_k}
 \longrightarrow \log G-\log F                         \tag{6.0d}
\end{equation}
uniformly on smaller discs.  Differentiating at zero shows every coefficient $b_j(L_{n_k})$
converges.  Limit on right of (6.0d) is corresponding polynomial $P_b$, so $G=\Phi_A(b)$ by
analytic continuation.  This proves cluster identity (1.U).

Continuity of $\Phi_A$ is immediate.  Its inverse on image is continuous because coefficients
are recovered by
\[
 b_j=\frac1{j!}\frac{d^j}{dz^j}\bigg|_{z=0}
       \log\frac{\Phi_A(b)(z)}{F(z)}.
\]
Thus $\Phi_A$ is topological embedding.  Exponential factor is zero-free, proving common-divisor
claim.  Cancelling $F$ gives (1.T), and cocycle law is additive identity for polynomials.  If jet
cluster set has nonempty interior, its difference set contains a neighborhood of zero; finite
sums of that neighborhood fill $\R^{p-1}$, proving full generated character group.  This
completes proof of Theorem~\ref{thm:universal-law}.

At $p=2$, (1.0) reads $b_1(L)=-\FP_{s=1}\zeta_L(s)$.  We now compute and realize this coordinate
for arithmetic operators.

Expanding finite Euler factor in (1.4) at $s=1$ gives
\[
 \prod_{p\in P}(1-p^{-s})
 =a(P)+a(P)\left(\sum_{p\in P}\frac{\log p}{p-1}\right)(s-1)
 +O((s-1)^2).
\]
Since $\zeta(s)=(s-1)^{-1}+\gamma+O(s-1)$, this proves (1.5).

Let
\[
 H_P(N)=\sum_{\substack{m\leq N\\(m,M_P)=1}}\frac1m.
\]
Equation (1.2) says
\begin{equation}
 H_P(N)=a(P)\log N+\gamma(P)+o(1).                         \tag{6.1}
\end{equation}
For cutoff product in (1.6), write
\begin{align}
 N^{a(P)z}\prod_{\substack{m\leq N\\(m,M_P)=1}}
 \left(1-\frac zm\right)
 &=\exp\{z(a(P)\log N-H_P(N))\}                            \notag\\
 &\quad\times
 \prod_{\substack{m\leq N\\(m,M_P)=1}}
 \left(1-\frac zm\right)e^{z/m}.                          \tag{6.2}
\end{align}
First factor tends locally uniformly to $e^{-\gamma(P)z}$ by (6.1).  Second converges locally
uniformly because $\sum m^{-2}<\infty$; by definition its limit is
$\det{}_2(I-zQ_P)$.  This proves (1.6) and (1.7).

To obtain (1.8), decompose coprime integers into reduced residue classes modulo $M=M_P$.
For $K\geq0$,
\begin{equation}
 \prod_{k=0}^{K}\left(1-\frac{z}{r+kM}\right)
 =\frac{\Gamma(K+1+(r-z)/M)}{\Gamma((r-z)/M)}
  \frac{\Gamma(r/M)}{\Gamma(K+1+r/M)}.                   \tag{6.3}
\end{equation}
Uniform Gamma-ratio asymptotic on compact $z$-sets contributes $K^{-z/M}$.  Multiplication over
$\varphi(M)$ reduced residue classes, followed by $N^{a(P)z}$ correction and
$a(P)=\varphi(M)/M$, gives (1.8).

Same calculation identifies zeta determinant.  On a reduced residue class,
\[
 \sum_{k\geq0}(r+kM-z)^{-s}=M^{-s}\zeta_H\left(s,\frac{r-z}{M}\right),
\]
where $\zeta_H$ is Hurwitz zeta function.  Classical values
\begin{equation}
 \zeta_H(0,w)=\frac12-w,\qquad
 \zeta_H'(0,w)=\log\Gamma(w)-\frac12\log(2\pi)            \tag{6.4}
\end{equation}
show that
\[
 -\frac{d}{ds}\bigg|_{s=0}
 \left(\Tr((D_P-z)^{-s})-\Tr(D_P^{-s})\right)
 =\log\Delta_P(z).
\]
This is (1.9), initially off positive spectrum and then everywhere by entire extension.

For prime excision, multiples of $p$ inside $P$-coprime integers are in bijection with all
$P$-coprime integers after division by $p$.  Therefore
\begin{equation}
 \det{}_2(I-zQ_{P\cup\{p\}})
 =\frac{\det{}_2(I-zQ_P)}{\det{}_2(I-(z/p)Q_P)}.           \tag{6.5}
\end{equation}
Lemma~\ref{lem:update}, or direct multiplication, also gives
\begin{equation}
 \gamma(P\cup\{p\})
 =\left(1-\frac1p\right)\gamma(P)+\frac{a(P)\log p}{p}.  \tag{6.6}
\end{equation}
Substitution of (6.5) and (6.6) into (1.7) proves excision law (1.10), completing proof of
Theorem~\ref{thm:determinant}.

\section{Cofinal determinant clusters and no descent}
\label{sec:nodescent}

Embed every $H_P$ in $H=\ell^2(\N_{>0})$ by extending $Q_P$ by zero.  Along an exhaustive
sequence,
\begin{equation}
 \norm{Q_P-|e_1\rangle\langle e_1|}_2^2
 =\sum_{\substack{m\geq2\\(m,M_P)=1}}\frac1{m^2}
 \longrightarrow0                                         \tag{7.1}
\end{equation}
by dominated convergence: every $m>1$ has a prime divisor eventually inserted.
Continuity of $\det{}_2$ in Hilbert--Schmidt norm~\cite{Simon2005} yields
\begin{equation}
 \det{}_2(I-zQ_P)\longrightarrow(1-z)e^z                \tag{7.2}
\end{equation}
locally uniformly.  Combining (7.2) with (1.7) proves forward assertion in
Theorem~\ref{thm:clusters}.

Conversely,
\begin{equation}
 \Delta_P'(0)=-\gamma(P),                                 \tag{7.3}
\end{equation}
because first derivative of $2$-regularized determinant is zero.  Local uniform convergence
preserves derivatives at origin, so a convergent determinant subsequence forces corresponding
$\gamma(P)$ to converge to a finite value.  Equations (7.2) and (7.3) identify its limit uniquely
as (1.11).  Formula (1.13) is immediate.

Give triples
\begin{equation}
 \mathfrak p(P)=\bigl(Q_P,a(P),\Div\Delta_P\bigr)           \tag{7.4}
\end{equation}
product topology of Hilbert--Schmidt norm, Euclidean topology, and local convergence of divisors.
Along every prime ordering,
\begin{equation}
 \mathfrak p(P_n^\sigma)\longrightarrow
 \bigl(|e_1\rangle\langle e_1|,0,[1]\bigr).               \tag{7.5}
\end{equation}
Residue limit follows from $a(P_n^\sigma)\leq a(E_k)\to0$ after $E_k$ has been included.
Local divisor convergence follows because each bounded set contains only finitely many positive
integer zeros, and every integer $m>1$ is eventually removed by one of its prime divisors.

\begin{corollary}[No descent from primary determinant data]
\label{cor:nodescent}
There is no continuous rule on closure of data \emph{(7.4)} that recovers cofinal
residue-corrected determinant.  Fibre over single primary limit \emph{(7.5)} contains all germs
\[
 e^{(1-c)z}(1-z),\qquad c\in[e^{-\gamma},\infty),
\]
and its relative-torsion group completion is
\begin{equation}
 \{e^{rz}:r\in\R\}.                                      \tag{7.6}
\end{equation}
\end{corollary}

\begin{proof}
Singleton case of Theorem~\ref{thm:classification} realizes every finite
$c\geq e^{-\gamma}$ while (7.5) stays fixed.  Theorem~\ref{thm:clusters} gives displayed fibre.
Pairwise quotients have slope $d-c$; all real slopes occur by choosing and reversing two points
of half-line.
\end{proof}

Finite part $\gamma(P)=\FP_{s=1}Z_P(s)$ is therefore not an optional normalization.  It is exact
secondary coordinate required for Hausdorff descent of determinant character.

\subsection{Full residue-jet sharpness by graded spectral suspension}

We prove Theorem~\ref{thm:full-jet}.  Fix $p\geq2$.  For $1\leq j<p$ define spectral
suspension on $P$-coprime integers by
\[
 D_P^{[j]}e_m=m^{1/j}e_m,\qquad
 L_{j,r,P}=\frac1{r+1}D_P^{[j]},\quad 0\leq r\leq j.
\]
Its inverse belongs to $\mathcal S_p$ because $p/j>1$, and
\[
 \zeta_{L_{j,r,P}}(s)=(r+1)^s Z_P(s/j).
\]
Let
\begin{equation}
 w_{j,r}=(-1)^{j-r}\binom jr,\qquad
 V_j(s)=\sum_{r=0}^j w_{j,r}(r+1)^s.                    \tag{7.7}
\end{equation}
Finite-difference identity gives
\begin{equation}
 V_j(k)=0\quad(0\leq k<j),\qquad V_j(j)=j!.             \tag{7.8}
\end{equation}

Take $p-1$ independent finite prime sets $\mathbf P=(P_1,\ldots,P_{p-1})$.  Graded virtual
spectral zeta function and normalized determinant character are
\begin{align}
 \mathfrak Z_{\mathbf P}(s)
   &=\sum_{j=1}^{p-1}\sum_{r=0}^j w_{j,r}\zeta_{L_{j,r,P_j}}(s)
     =\sum_{j=1}^{p-1}V_j(s)Z_{P_j}(s/j),                \tag{7.9}\\
 \mathfrak D_{\mathbf P}(z)
   &=\prod_{j=1}^{p-1}\prod_{r=0}^j
     \left(\frac{\det_\zeta(L_{j,r,P_j}-z)}
                  {\det_\zeta L_{j,r,P_j}}\right)^{w_{j,r}}. \notag
\end{align}
This virtual notation is an ordinary finite superdeterminant.  Put $|w_{j,r}|$ copies of
$L_{j,r,P_j}$ in even degree when $w_{j,r}>0$ and in odd degree when $w_{j,r}<0$.  Then
$\mathfrak Z_{\mathbf P}$ is supertrace spectral zeta function and
$\mathfrak D_{\mathbf P}$ is normalized section of graded determinant line.  Negative weights
mean inverse determinant sections; meromorphic character appears only because common even/odd
divisors have not yet been cancelled.  Thus construction belongs to standard determinant-line
and analytic-torsion formalism, not formal subtraction of positive operators.  Formula (1.0)
extends additively to corrected graded jet
$\mathfrak b=(\mathfrak b_1,\ldots,\mathfrak b_{p-1})$.

At $s=k$, summands with $j>k$ vanish by (7.8), while those with $j<k$ are regular.  At $j=k$,
\[
 Z_{P_k}((k+t)/k)=\frac{k a(P_k)}t+\gamma(P_k)+O(t).
\]
Therefore exact corrected coordinate is
\begin{equation}
\begin{split}
 \mathfrak b_k(\mathbf P)
  ={}&-\frac{V_k(k)}k\gamma(P_k)
      -\bigl(V_k'(k)+H_{k-1}V_k(k)\bigr)a(P_k)\\
    &-\frac1k\sum_{j<k}V_j(k)Z_{P_j}(k/j).
\end{split}
\tag{7.10}
\end{equation}
Along any coordinatewise exhaustion, $a(P_k)\to0$ and
$Z_{P_j}(t)\to1$ for every $t>1$.  Put
\begin{equation}
 d_k=-\frac{V_k(k)}k=-(k-1)!,\qquad
 C_k=-\frac1k\sum_{j<k}V_j(k).                          \tag{7.11}
\end{equation}
Then
\begin{equation}
 \mathfrak b_k(\mathbf P)=C_k+d_k\gamma(P_k)+o(1).       \tag{7.12}
\end{equation}
Thus finite-difference grading cancels every lower coordinate exactly and isolates one Mertens
finite part on diagonal.

Primary data converge.  Indeed, for each constituent,
\begin{equation}
 \left\|L_{j,r,P}^{-1}-(r+1)|e_1\rangle\langle e_1|\right\|_p^p
 =(r+1)^p\sum_{\substack{m\geq2\\(m,M_P)=1}}m^{-p/j}\longrightarrow0. \tag{7.13}
\end{equation}
Only $j=k$ term of (7.9) has pole at integer $k$, with residue
$kV_k(k)a(P_k)\to0$.  Divisor convergence is locally eventually exact away from surviving
rank-one zero.  Indeed zeros of constituent determinant are
\[
 z=(r+1)^{-1}m^{1/j},\qquad (m,M_P)=1.
\]
On a fixed compact set only finitely many such integers $m$ occur.  Every $m>1$ has a prime
divisor and is therefore eventually removed by any exhaustion, while $m=1$ remains.  Hence
every constituent divisor converges locally to fixed divisor at $z=(r+1)^{-1}$, and finite
graded virtual divisor converges as well.

It remains to determine generic secondary spectrum.  Let $\Omega_p$ be Polish space of all
orderings of colored primes $\{1,\ldots,p-1\}\times\PP$; let $P_{j,n}$ contain primes of color
$j$ seen by time $n$.  For any basic open rectangle in
$[e^{-\gamma},\infty]^{p-1}$ and any finite ordered prefix, extend prefix one color at a time.
Rational finite-prefix extension, obtained by starting construction of Section~5 after given
finite set, places $\gamma(P_{j,n})$ in prescribed coordinate neighborhood without changing
other colors.  Hence event of hitting rectangle after time $N$ is open dense.  Intersecting over
countable rectangle basis and $N$ proves that on comeagre set of orderings,
\begin{equation}
 \operatorname{Clust}_{n\to\infty}
   \bigl(\gamma(P_{1,n}),\ldots,\gamma(P_{p-1,n})\bigr)
 =[e^{-\gamma},\infty]^{p-1}.                            \tag{7.14}
\end{equation}
Equations (7.12) and (7.14) show finite virtual-jet cluster set is product
\[
 \prod_{k=1}^{p-1}\bigl(C_k+d_k[e^{-\gamma},\infty)\bigr),
\]
with nonempty interior in $\R^{p-1}$.  Theorem~\ref{thm:universal-law} applies to finite graded
sums and differences after common virtual divisor cancellation.  Its character-group conclusion
gives (1.F).  Varying only $k$th scalar cluster coordinate changes torsion by
$\exp(d_k(c-c')z^k)$ and leaves all other coordinates fixed; hence no $k$th jet coordinate can
be deleted.  This proves Theorem~\ref{thm:full-jet}.

\begin{corollary}[Sharp no-descent from primary determinant data]
\label{cor:full-no-descent}
For every $p\geq2$, no rule continuous at common limiting primary datum can recover virtual
zeta-determinant character from constituent $\mathcal S_p$ limits, integer zeta residues, and
local limiting divisors alone.  More sharply, after adjoining any proper subset of corrected
coordinates, omitted coordinate still carries full one-parameter torsion group
$\{\exp(vz^k):v\in\R\}$.
\end{corollary}

\begin{proof}
Along family of Theorem~\ref{thm:full-jet}, all listed primary data converge independently of
colored-prime ordering.  A continuous recovery rule would therefore force unique determinant
limit.  Equation (7.14) and (7.12) instead give determinant cluster set with nonempty interior.
If only coordinate $k$ is omitted, hold all other scalar cluster coordinates fixed and vary
$\gamma(P_k)$ over its half-line.  Pairwise differences fill $\R$, since $d_k\neq0$, and
relative torsions are exactly $\exp(vz^k)$.  Thus no proper sub-jet repairs descent.
\end{proof}

\section{Number fields and universal residue cancellation}
\label{sec:numberfields}

Let $K$ be a number field, let $\mathcal O_K$ be its ring of integers, and let
$\mathcal P_K$ be set of nonzero prime ideals.  Write
\begin{equation}
 \zeta_K(s)=\frac{\rho_K}{s-1}+c_K+O(s-1),\qquad
 \gamma_K=\frac{c_K}{\rho_K}.                            \tag{8.1}
\end{equation}
Fix once and for all a reference ordering
$\mathfrak p_1\prec\mathfrak p_2\prec\cdots$ which refines nondecreasing norm; ties are resolved
arbitrarily.  Put $E_r^K=\{\mathfrak p_1,\ldots,\mathfrak p_r\}$.  For finite
$P\subset\mathcal P_K$, define
\begin{align}
 a_K(P)&=\prod_{\mathfrak p\in P}\left(1-\frac1{N\mathfrak p}\right),
 &r_K(P)&=\rho_Ka_K(P),                                  \tag{8.2}\\
 b_K(P)&=\gamma_K+\sum_{\mathfrak p\in P}
                   \frac{\log N\mathfrak p}{N\mathfrak p-1},
 &\kappa_K(P)&=r_K(P)b_K(P).                             \tag{8.3}
\end{align}
Kandhil--Lunia introduced normalized generalized Euler--Kronecker constants; for finite $P$ their
constant is $a_K(P)b_K(P)$~\cite{KandhilLunia2024}.  Our $\kappa_K(P)$ is unnormalized finite
part relevant to zeta determinant.

\begin{proposition}[Pole data and universal endpoint]
\label{prop:K-pole}
For
\begin{equation}
 Z_{K,P}(s)=\zeta_K(s)\prod_{\mathfrak p\in P}
                    (1-N\mathfrak p^{-s})                \tag{8.4}
\end{equation}
one has
\begin{equation}
 \Res_{s=1}Z_{K,P}(s)=r_K(P),\qquad
 \FP_{s=1}Z_{K,P}(s)=\kappa_K(P).                        \tag{8.5}
\end{equation}
Moreover,
\begin{equation}
 \sum_{\substack{N\mathfrak a\leq X\\(\mathfrak a,P)=1}}
 \frac1{N\mathfrak a}
 =r_K(P)\log X+\kappa_K(P)+o(1),                         \tag{8.6}
\end{equation}
and
\begin{equation}
 r_K(E_r^K)b_K(E_r^K)=\kappa_K(E_r^K)\longrightarrow e^{-\gamma}.
                                                                    \tag{8.7}
\end{equation}
In particular, endpoint is independent of $K$.
\end{proposition}

\begin{proof}
Expanding each Euler factor at $1$ gives
\[
 \prod_{\mathfrak p\in P}(1-N\mathfrak p^{-s})
 =a_K(P)\left(1+(s-1)\sum_{\mathfrak p\in P}
 \frac{\log N\mathfrak p}{N\mathfrak p-1}+O((s-1)^2)\right).
\]
Together with (8.1), this proves (8.5).  Partial summation from ideal-counting asymptotic, or
finite-set formula of Kandhil--Lunia, gives (8.6).

Number-field Mertens theorems give, with $E^K(X)=\{\mathfrak p:N\mathfrak p\leq X\}$,
\begin{equation}
 a_K(E^K(X))\sim\frac{e^{-\gamma}}{\rho_K\log X},\qquad
 \sum_{N\mathfrak p\leq X}\frac{\log N\mathfrak p}{N\mathfrak p-1}
 =\log X+C_K+o_K(1)                                     \tag{8.8}
\end{equation}
see Rosen and unconditional explicit form of Garcia--Lee
\cite{Rosen1999,GarciaLee2022}.  Multiplying (8.8) by $\rho_K$ proves (8.7).  At most
$[K:\mathbb Q]$ prime ideals have any fixed norm, so passing through a tie changes product by
$1+O_K(X^{-1})$ and weighted sum by $O_K(\log X/X)$.  Same limit holds for every $E_r^K$.
\end{proof}

Next lemma is number-field version of relative extremality.  Repeated norms require an explicit
termination argument.

\begin{lemma}[Prime-ideal relative extremality]
\label{lem:K-extremal}
There is $k_0(K)$ such that, for $k\geq k_0(K)$ and every finite $P$ containing $E_k^K$,
\begin{equation}
 \kappa_K(P)\geq\inf_{r\geq k}\kappa_K(E_r^K).          \tag{8.9}
\end{equation}
Consequently every increasing exhaustion of $\mathcal P_K$ satisfies
\begin{equation}
 \liminf_n\kappa_K(P_n)\geq e^{-\gamma}.               \tag{8.10}
\end{equation}
\end{lemma}

\begin{proof}
Choose $k_0$ so that $b_K(E_{k_0}^K)>0$.  Start with $Q\supset E_k^K$.  If $Q$ is not an
initial segment, let $\mathfrak t$ be its largest member in reference order, set
$Q'=Q\setminus\{\mathfrak t\}$, and write $A=b_K(Q')>0$.  If
$\kappa_K(Q')<\kappa_K(Q)$, replace $Q$ by $Q'$.  Otherwise let $\mathfrak u$ be least missing
prime ideal and replace $Q$ by $Q'\cup\{\mathfrak u\}$.

For a prime ideal of norm $v$ not in $Q'$,
\[
 \kappa_K(Q'\cup\{\mathfrak p\})=\kappa_K(Q')f_A(v),\qquad
 f_A(v)=1-\frac1v+\frac{\log v}{vA}.
\]
Function $f_A$ increases up to $e^{A+1}$ and decreases afterwards.  In replacement case
$f_A(N\mathfrak t)\leq1$, hence $N\mathfrak t\leq e^A$.  Since
$\mathfrak u\prec\mathfrak t$, we have $N\mathfrak u\leq N\mathfrak t$ and thus
$f_A(N\mathfrak u)\leq f_A(N\mathfrak t)$.  Value never increases.

Let $i(\mathfrak p)$ be reference index.  At each deletion or replacement,
$\sum_{\mathfrak p\in Q}i(\mathfrak p)$ strictly decreases.  Thus process terminates, including
when $N\mathfrak u=N\mathfrak t$, at an initial segment $E_r^K$.  Every move preserves $E_k^K$,
so $r\geq k$ and (8.9) follows.  Apply (8.9) after fixed $E_k^K$ enters exhaustion, then let
$k\to\infty$ and use (8.7) to prove (8.10).
\end{proof}

Exact local dynamics is
\begin{equation}
 \kappa_K(P\cup\{\mathfrak p\})-\kappa_K(P)
 =\frac{r_K(P)}{N\mathfrak p}
   \bigl(\log N\mathfrak p-b_K(P)\bigr).                \tag{8.11}
\end{equation}
Therefore, if $\kappa_K(P)\leq M$ and $q=N\mathfrak p$, then
\begin{equation}
 |\kappa_K(P\cup\{\mathfrak p\})-\kappa_K(P)|
 \leq\frac{M+\rho_K\log q}{q}.                         \tag{8.12}
\end{equation}
This is vanishing mesh because any enumeration of $\mathcal P_K$ has $N\mathfrak p_n\to\infty$.

For $y>1$ and $u\geq0$, put
\[
 B_K(y,u)=\{\mathfrak p:y<N\mathfrak p\leq e^uy\}.
\]
Prime ideal theorem and (8.8) imply, for fixed finite $P$,
\begin{align}
 \prod_{\mathfrak p\in B_K(y,u)}(1-N\mathfrak p^{-1})&\longrightarrow1,
                                                                    \tag{8.13}\\
 \sum_{\mathfrak p\in B_K(y,u)}
       \frac{\log N\mathfrak p}{N\mathfrak p-1}&\longrightarrow u,
                                                                    \tag{8.14}\\
 \kappa_K(P\cup B_K(y,u))&\longrightarrow
             \kappa_K(P)+r_K(P)u.                                 \tag{8.15}
\end{align}
Taking $\log y>b_K(P)+u+1$ makes every insertion in block increase value; (8.12) makes its mesh
arbitrarily small.

Conversely, let $\mathfrak q=\mathfrak p_k$ be least missing prime ideal.  Then
$a_K(P)\leq a_K(E_{k-1}^K)$ and (8.7) gives
\begin{equation}
 \rho_Ka_K(E_{k-1}^K)\log N\mathfrak p_k\longrightarrow e^{-\gamma}.
                                                                    \tag{8.16}
\end{equation}
Hence, for fixed $\alpha>e^{-\gamma}$ and all large $k$, every least-missing insertion made while
$\kappa_K(P)\geq\alpha$ decreases value.  For $\alpha=e^{-\gamma}$ same holds outside any fixed
neighborhood of endpoint.

\subsection{Abstract Mertens-system universality}

Concrete number-field proof exposes exact general mechanism.  Let $\mathcal A$ be countably
infinite atom set with proper height $q:\mathcal A\to(1,\infty)$, meaning each bounded sublevel is
finite.  Fix $\rho>0$ and $\gamma_0\in\R$.  For finite $P\subset\mathcal A$, put
\begin{equation}
 a(P)=\prod_{v\in P}(1-q(v)^{-1}),\quad
 b(P)=\gamma_0+\sum_{v\in P}\frac{\log q(v)}{q(v)-1},\quad
 \kappa(P)=\rho a(P)b(P).                                \tag{8.17}
\end{equation}
For $x>1$, write $\mathcal A(x)=\{v:q(v)\leq x\}$ and suppose
\begin{equation}
 \rho a(\mathcal A(x))\log x\longrightarrow c_*>0,\qquad
 \sum_{q(v)\leq x}\frac{\log q(v)}{q(v)-1}-\log x
 \longrightarrow C.                                    \tag{8.18}
\end{equation}

\begin{theorem}[Mertens-system universality]
\label{thm:mertens-system}
Under (8.18), cluster sets of $\kappa(P_n)$ along arbitrary orderings of all atoms are exactly
nonempty closed intervals in $[c_*,\infty]$.  Every interval occurs.  Every interval fibre is
dense in ordering space, full interval $[c_*,\infty]$ is comeagre, and cluster-spectrum map is
nowhere continuous.
\end{theorem}

\begin{proof}
Order atoms by nondecreasing height, with arbitrary tie-breaking.  Second limit in (8.18) implies
total jump across every equal-height tie tends to zero; first then implies, also through partial
ties,
\begin{equation}
 \kappa(E_n)\to c_*,\qquad
 \rho a(E_{n-1})\log q(v_n)\to c_*.                      \tag{8.19}
\end{equation}
Indeed $b(\mathcal A(x))/\log x\to1$; convergence of second expression in (8.18) forces
$\log(q(v_{n+1})/q(v_n))\to0$ and makes every tie negligible.

One-atom update is
\[
 \kappa(P\cup\{v\})-\kappa(P)
 =\frac{\rho a(P)}{q(v)}(\log q(v)-b(P)).
\]
Relative extremal descent of Lemma~\ref{lem:K-extremal} uses only this formula, height order, and
eventual positivity of $b(E_n)$; its reference-index potential also handles ties.  Thus every
cofinal liminf is at least $c_*$.  On bounded windows, jump is at most
$(M+\rho\log q(v))/q(v)$ and hence tends to zero.

For fixed $u\geq0$, divide first limit in (8.18) at $e^uy$ and $y$, and subtract second limit.
This gives
\begin{align*}
 \prod_{y<q(v)\leq e^uy}(1-q(v)^{-1})&\to1,\\
 \sum_{y<q(v)\leq e^uy}\frac{\log q(v)}{q(v)-1}&\to u.
\end{align*}
These are remote upward control; (8.19) is downward control.  Crossing proof and alternating
block construction used above now apply word for word with $c_*$ in place of $e^{-\gamma}$.
Finite-prefix extension and Baire proof of Section~\ref{sec:generic} use only same controls, so
density, genericity, and nowhere continuity follow as well.
\end{proof}

\begin{proof}[Proof of Theorem~\ref{thm:numberfield-main}]
Let $(\mathfrak q_n)$ be any permutation of $\mathcal P_K$.  Closedness of cluster set is
automatic in compact space $[0,\infty]$, and (8.10) gives lower bound.  If two finite cluster
values surround $t$, sequence crosses $t$ infinitely often.  On bounded crossing window,
(8.12) makes overshoot tend to zero; thus $t$ is also cluster value.  Same argument at an
unbounded end proves cluster set is interval in $[e^{-\gamma},\infty]$.

Conversely fix $e^{-\gamma}\leq\alpha\leq\beta\leq\infty$.  Alternate two finite operations.
Use remote block (8.15), with $u$ chosen from affine limit
$\kappa_K(P)+r_K(P)u$, to cross upward target $\beta$ with prescribed vanishing mesh; if
$\beta=\infty$, use targets tending to infinity.  Then insert least missing prime ideals until
crossing $\alpha$ downward; if $\alpha=e^{-\gamma}$, continue to an initial segment close to
endpoint.  At every downward stage consume at least one new least missing ideal.  Equations
(8.12), (8.15), and (8.16) give errors tending to zero.  Every prime ideal is eventually
inserted, all finite cluster values lie in $[\alpha,\beta]$, and vanishing-mesh crossings make
every point of interval a cluster value.  Same construction with coincident targets gives
singletons, using consume-then-restore modification of Section~5; targets tending to infinity
without downward sweeps give $\{\infty\}$.  This proves cluster classification.

For determinant lift, let $H_{K,P}$ have orthonormal basis indexed by nonzero integral ideals
coprime to $P$ and put
\[
 D_{K,P}e_{\mathfrak a}=N\mathfrak a\,e_{\mathfrak a},\qquad
 Q_{K,P}=D_{K,P}^{-1}.
\]
Then $Q_{K,P}$ is Hilbert--Schmidt because $Z_{K,P}(2)<\infty$, and (8.4) is its spectral zeta
function.  From (8.6), same factorization as (6.2) gives locally uniformly on $\C$
\begin{align}
 \Delta_{K,P}(z)
 &:=\lim_{X\to\infty}X^{r_K(P)z}
 \prod_{\substack{N\mathfrak a\leq X\\(\mathfrak a,P)=1}}
       \left(1-\frac z{N\mathfrak a}\right)                       \tag{8.20}\\
 &=e^{-\kappa_K(P)z}\det{}_2(I-zQ_{K,P}).                         \tag{8.21}
\end{align}

This is exactly relative zeta determinant, not merely a cutoff analogy.  For $|z|<1$, set
\[
 \Xi_{K,P}(s,z)=\Tr\bigl((D_{K,P}-z)^{-s}-D_{K,P}^{-s}\bigr).
\]
Binomial expansion gives meromorphic identity near $s=0$
\begin{equation}
 \Xi_{K,P}(s,z)=\sum_{m\geq1}\frac{(s)_m}{m!}z^mZ_{K,P}(s+m).
                                                                    \tag{8.22}
\end{equation}
Since $sZ_{K,P}(1+s)=r_K(P)+\kappa_K(P)s+O(s^2)$ and
$(s)_m'|_{s=0}=(m-1)!$ for $m\geq2$,
\[
 -\Xi_{K,P}'(0,z)=-\kappa_K(P)z-
     \sum_{m\geq2}\frac{Z_{K,P}(m)}m z^m=\log\Delta_{K,P}(z).
\]
Thus
\begin{equation}
 \Delta_{K,P}(z)=\frac{\det_\zeta(D_{K,P}-z)}{\det_\zeta D_{K,P}},
                                                                    \tag{8.23}
\end{equation}
first for $|z|<1$ and then by entire continuation.  If $\mathfrak p\notin P$ and
$q=N\mathfrak p$, ideals divisible by $\mathfrak p$ are $\mathfrak p\mathfrak a$ uniquely.
Using
\[
 \kappa_K(P\cup\{\mathfrak p\})=\left(1-\frac1q\right)\kappa_K(P)
 +\frac{r_K(P)\log q}{q}
\]
in (8.21) yields exact prime-ideal excision
\begin{equation}
 \boxed{\displaystyle
 \Delta_{K,P\cup\{\mathfrak p\}}(z)
 =q^{-r_K(P)z/q}\frac{\Delta_{K,P}(z)}{\Delta_{K,P}(z/q)}.}       \tag{8.24}
\end{equation}
Equation (8.24) is identity of meromorphic quotients with removable singularities; globally
equivalent entire identity is
\[
 q^{r_K(P)z/q}\Delta_{K,P\cup\{\mathfrak p\}}(z)
 \Delta_{K,P}(z/q)=\Delta_{K,P}(z).
\]

Embed all operators in $\ell^2$ of nonzero integral ideals.  Along every exhaustion,
\[
 Q_{K,P}\longrightarrow|e_{\mathcal O_K}\rangle\langle e_{\mathcal O_K}|
 \quad\hbox{in Hilbert--Schmidt norm},
\]
because every nonunit ideal has a prime-ideal divisor.  Also $r_K(P)\to0$ by (8.8), while
divisors converge locally to $[1]$.  Continuity of $\det_2$ now gives, whenever
$\kappa_K(P_{n_j})\to c<\infty$,
\[
 \Delta_{K,P_{n_j}}(z)\longrightarrow e^{(1-c)z}(1-z).
\]
Derivative at zero equals $-\kappa_K(P)$, proving converse.  Quotients are $e^{(d-c)z}$.
Ordering-independence of primary data and all remaining claims follow.
\end{proof}

\section{Generic maximal torsion and topological instability}
\label{sec:generic}

Topology on $\mathfrak S_K$ has basic cylinders obtained by fixing a finite ordered prefix.
It is Polish: identify an ordering with a bijection $\N_{>0}\to\mathcal P_K$ inside product of
countable discrete spaces.  Write $C_K(\sigma)$ for cluster set of $\kappa_K(P_n^\sigma)$.

\begin{lemma}[Finite-prefix extension]
\label{lem:prefix}
Every finite ordered prefix extends to an ordering with any prescribed cluster interval
$I\subset[e^{-\gamma},\infty]$.
\end{lemma}

\begin{proof}
Start interval-realization construction in proof of Theorem~\ref{thm:numberfield-main} from set
already used by prefix.  If necessary, first insert least missing prime ideals through a large
reference cutoff.  This finite initialization does not affect cluster set.  Remote blocks can
then be chosen beyond all used norms, while every downward stage consumes new least missing
ideal.  Result is permutation extending prefix and has prescribed cluster interval.
\end{proof}

For open interval $J\subset(e^{-\gamma},\infty)$ and $N\geq1$, let
\[
 U(J,N)=\{\sigma:\kappa_K(P_n^\sigma)\in J\text{ for some }n\geq N\}.
\]
This set is open because witness uses finite prefix, and dense by
Lemma~\ref{lem:prefix}, choosing any interval meeting $J$.  Likewise
\[
 U_\infty(M,N)=\{\sigma:\kappa_K(P_n^\sigma)>M\text{ for some }n\geq N\}
\]
is open dense.  Intersect $U(J,N)$ over countable rational basis $J$, all $N$, and intersect
$U_\infty(M,N)$ over positive integers $M,N$.  Baire theorem shows resulting set is comeagre.
Every ordering in it has every finite $c\geq e^{-\gamma}$ and $\infty$ as cluster values; lower
bound (8.10) excludes others.  Thus full-spectrum fibre is comeagre.  Lemma~\ref{lem:prefix}
makes every interval fibre dense, while each proper fibre lies in meagre complement of full one.

Equip $[e^{-\gamma},\infty]$ with any compatible compact metric and its nonempty-compact
hyperspace with Hausdorff metric.  Every cylinder contains, by Lemma~\ref{lem:prefix}, orderings
with cluster sets $\{e^{-\gamma}\}$ and $\{\infty\}$.  These two hyperspace points have positive
distance.  Triangle inequality prevents both from lying in arbitrarily small neighborhood of
single $C_K(\sigma)$.  Cluster-spectrum map is therefore continuous nowhere.  Determinant
claim follows from Theorem~\ref{thm:numberfield-main}: every finite cluster parameter $c$ gives
$e^{(1-c)z}(1-z)$ and pairwise quotients give every slope.  This proves
Theorem~\ref{thm:generic}.

\section{Quillen metrics, Wodzicki residues, and analytic torsion}
\label{sec:geometry}

We record precise relation to geometric determinant literature and limit of that analogy.  If
$A$ is a positive classical pseudodifferential operator of order $-d$ on a closed $d$-manifold,
then canonical trace of holomorphic family $A^s$ has
\begin{equation}
 \Res_{s=1}\Tr(A^s)=\frac1d\operatorname{res}_W(A),        \tag{10.1}
\end{equation}
where $\operatorname{res}_W$ is Wodzicki residue~\cite{PaychaScott2007}.  Equation (10.1) is classical
pseudodifferential counterpart of $\Res_{s=1}Z_P(s)=a(P)$.  We do not claim diagonal arithmetic
operator $Q_P$ is a classical pseudodifferential operator.  Common structure is a simple critical
zeta pole and a residue-controlled logarithmic counterterm.

More generally, for positive elliptic $L$ of order $m$,
\[
 \Res_{s=j}\zeta_L(s)=\frac1m\operatorname{res}_W(L^{-j}).
\]
Consequently universal coordinate (1.0) splits canonically into nonlocal finite part and local
Wodzicki correction:
\[
 b_j(L)=-\frac1j\FP_{s=j}\zeta_L(s)
        -\frac{H_{j-1}}{jm}\operatorname{res}_W(L^{-j}).
\]
Theorem~\ref{thm:universal-law} says entire vector of these mixed coordinates, rather than
determinant line or residue vector alone, is complete cofinal datum.

This also separates two polynomial-gauge statements.  Dang's elliptic theorem identifies
renormalized determinant trivializations obtainable from local counterterms; their ambiguity is
torsor under local polynomial functionals~\cite{Dang2022}.  Restrict such functional to any
finite-dimensional perturbation pencil whose inverse perturbation lies in $\mathcal S_p$: its
normalized degree-$<p$ part is a section of gauge group in
Theorem~\ref{thm:parametrized-completion}.  Our completion theorem identifies ambient
Schatten-topological gauge group and proves its full realization by raw determinant limits; it
does not assert that every ambient direction is microlocally local.  In zeta lift above,
Wodzicki term is local part and finite part is generally nonlocal part of same fibre coordinate.

Paycha--Rosenberg prove regularized first Chern form of an infinite-dimensional bundle differs
from Bismut--Freed curvature by Wodzicki-residue terms~\cite{PaychaRosenberg2000}.  Scott residue
determinant is local and multiplicative, in sharp contrast with nonlocal zeta determinant
\cite{Scott2005}.  Corollary~\ref{cor:nodescent} shows complementary arithmetic fact: even after
primary residue and limiting divisor are fixed, finite part at pole supplies a full real torsion
fibre.  This paper does not propose adding a Wodzicki residue to determinant line as new
construction; that general principle is established.  New content is universal cofinal
completeness of corrected jet and sharp arithmetic realization of its secondary spectrum.

\begin{corollary}[Gauge defect of determinant trivializations]
\label{cor:metric-connection}
Suppose two finite cofinal jet limits $b,c:B\to\R^{p-1}$ vary $C^1$ over parameter manifold $B$
and have same limiting $\det_p$ core in common local determinant-line trivialization.  At fixed
spectral parameter $z$, quotient section, ratio of its Quillen norms, and gauge transformation
of Bismut--Freed connection one-forms are
\[
 \mathcal T_{b,c}(z)=e^{P_b(z)-P_c(z)},\qquad
 \frac{\|\Phi_A(b)(z)\|_Q}{\|\Phi_A(c)(z)\|_Q}
   =e^{\operatorname{Re}(P_b(z)-P_c(z))},
\]
\[
 \omega_b(z)-\omega_c(z)=d_B(P_b(z)-P_c(z)).             \tag{10.M}
\]
Thus topological triviality and common divisor do not determine normalized determinant section
or its connection gauge.  Corrected residue jet determines exact transition character.  This
is a statement about two gauges of same metric and connection, not a claim that underlying
Quillen metric or Bismut--Freed connection changes.
\end{corollary}

\begin{proof}
First identity is (1.T).  Hermitian norm transforms by modulus of transition section, and
connection one-form transforms by its logarithmic differential.  Since exponential never
vanishes, logarithm is global polynomial $P_b-P_c$ in this trivialization.
\end{proof}

For fixed $P$, holomorphic Fredholm family $D_P-z$ has index zero.  Away from spectrum, (1.9) is
scalar of its relative determinant section against $D_P$ in spectral trivialization.  At
eigenvalues entire continuation records divisor of that section.  This is elementary diagonal
analogue of determinant-line/Quillen-metric relation for elliptic families
\cite{Quillen1985,BismutFreed1986,Scott1999,Bunke2002}.

There is also an exact analytic-torsion interpretation.  For real $z<1$, consider acyclic
two-term Hilbert complex
\[
 0\longrightarrow H_P\xrightarrow{D_P-z}H_P\longrightarrow0.
\]
Its zeta analytic torsion is $\det_\zeta(D_P-z)$ under standard two-term
convention~\cite{RaySinger1971}.  Hence
\begin{equation}
 \Delta_P(z)=\frac{T_{\rm an}(D_P-z)}{T_{\rm an}(D_P)},\qquad z<1. \tag{10.2}
\end{equation}
At two cofinal finite-part values $c,d$, relative analytic torsion ratio is $e^{(d-c)z}$.  For
complex parameters same scalar is a refined determinant torsion, retaining phase rather than
only Ray--Singer norm; compare~\cite{DaiFreed1994,BravermanKappeler2007}.  No geometric gluing
theorem is claimed for this diagonal model.

For finite graded family satisfying hypotheses of Theorem~\ref{thm:universal-law} degreewise,
same argument gives universal relative analytic-torsion character
\[
 \exp\!\left(\frac12\sum_q(-1)^q q
       \bigl(P_{b^{(q)}}(z)-P_{c^{(q)}}(z)\bigr)\right). \tag{10.A}
\]
This is relative determinant torsion after common degreewise divisors are cancelled.  Formula
(10.2) is $p=2$ two-term specialization.

\section{Schatten determinant completion and vanishing micro-spectra}
\label{sec:completion}

We prove parent law.  For finite-rank pencil $A:V\to\operatorname{Fin}(H)$, defining identity
of regularized determinant is
\begin{equation}
 D_A(v)=D_{p,A}(v)
 \exp\!\left(-\sum_{k=1}^{p-1}\frac1k\Tr(A(v)^k)\right).       \tag{11.1}
\end{equation}
Continuity of $A\mapsto D_{p,A}$ from
$\operatorname{Hom}(V,\mathcal S_p)$ to $\operatorname{Hol}_1(V)$ follows from standard
continuity estimate for $\det_p$, uniformly for $v$ in compact subsets of $V$
\cite{Simon2005}.

\subsection*{Necessity and topology of the fibre}

Suppose finite-rank $A_n\to A$ in $\operatorname{Hom}(V,\mathcal S_p)$ and
$D_{A_n}\to F$ compact-open.  By (11.1),
\[
 D_{A_n}=D_{p,A_n}e^{P_n},\qquad
 P_n(v)=-\sum_{k=1}^{p-1}\frac1k\Tr(A_n(v)^k)\in J_{<p}(V).
\]
Since $F(0)=D_{p,A}(0)=1$, choose ball $U$ about origin on which both are nonzero.  On every
smaller ball,
\[
 e^{P_n}=\frac{D_{A_n}}{D_{p,A_n}}\longrightarrow
 G:=\frac{F}{D_{p,A}}
\]
locally uniformly, and $G$ is nonzero.  Normalized holomorphic logarithm on $U$, with value zero
at origin, is unique.  It follows that $P_n\to\log G$ locally uniformly on $U$.  Equivalently,
all Taylor coefficients of $P_n$ converge.  Finite-dimensional space $J_{<p}(V)$ is closed, so
limit is some $P\in J_{<p}(V)$.  Identity theorem gives $F=e^P D_{p,A}$ on $V$.

Same argument proves inverse continuity in (1.D): from $(A_n,e^{P_n}D_{p,A_n})$ recover finite
jet $P_n$ by normalized logarithm of quotient near origin.  Thus $\Theta_{p,V}$ is
homeomorphism onto graph closure once sufficiency is proved.

\subsection*{Moment synthesis at vanishing Schatten cost}

We need elementary block.  Let $\ell\in V^*$, $c\in\C$, $1\leq k<p$, $N\geq1$, and choose
$t\in\C$ with $t^k=c/(kN)$.  If $\omega$ is primitive $k$th root of unity, define diagonal
pencil $\mathcal B_{k,c,\ell,N}$ with $N$ copies of each eigenvalue functional
\[
 t\ell,\ t\omega\ell,\ldots,t\omega^{k-1}\ell .
\]
Root-of-unity sum gives for every $q\geq1$
\begin{equation}
 \Tr\bigl(\mathcal B_{k,c,\ell,N}(v)^q\bigr)
 =
 \begin{cases}
  kNt^q\ell(v)^q,&k\mid q,\\
  0,&k\nmid q.
 \end{cases}                                               \tag{11.2}
\end{equation}
In particular all moments below $k$ vanish and $k$th moment is $c\ell^k$.  If
$\ell=\sum_i\ell_i e_i^*$ in fixed basis of $V$, Schatten cost of $i$th coefficient operator is
\begin{equation}
 \norm{\mathcal B_{k,c,\ell,N}(e_i)}_p^p
 =(kN)^{1-p/k}|c|^{p/k}|\ell_i|^p.                        \tag{11.3}
\end{equation}
It tends to zero because $k<p$.

Pure powers span $\operatorname{Sym}^k(V^*)$.  One way to see this is polarization identity:
coefficient of $t_1\cdots t_k$ in
$(t_1\ell_1+\cdots+t_k\ell_k)^k$ recovers symmetrization of
$\ell_1\otimes\cdots\otimes\ell_k$.  Hence every $Q\in\operatorname{Sym}^k(V^*)$ is finite sum
$Q=\sum_\nu c_\nu\ell_\nu^k$ and (11.2) realizes $Q$ as $k$th moment at arbitrarily small
$\mathcal S_p$ cost without changing lower moments.

Now let $A\in\operatorname{Hom}(V,\mathcal S_p)$ and
$P=\sum_{k=1}^{p-1}P_k$, $P_k\in\operatorname{Sym}^k(V^*)$.  Choose finite-dimensional
projections $E_n\uparrow I$ such that
\[
 A_n(v)=E_nA(v)E_n\longrightarrow A(v)
\]
in $\mathcal S_p$, uniformly on bounded subsets of $V$.  Construct on $E_n^\perp H$ a diagonal
pencil $B_n$ degree by degree.  Suppose blocks of degrees below $k$ already chosen and denote
their direct sum by $B_n^{<k}$.  Define residual homogeneous form
\[
 Q_{k,n}(v)
 =-kP_k(v)-\Tr(A_n(v)^k)-\Tr(B_n^{<k}(v)^k).
\]
Decompose $Q_{k,n}$ into pure $k$th powers and realize each summand by (11.2).  Later blocks have
zero moments of degrees at most $k$, so recursion preserves completed equations.  Multiplicities
$N$ may be chosen after residual coefficients are known.  By (11.3), choose them so total pencil
norm satisfies $\norm{B_n}_{\operatorname{Hom}(V,\mathcal S_p)}<1/n$.  At end,
\begin{equation}
 \Tr\bigl((A_n\oplus B_n)(v)^k\bigr)=-kP_k(v),
 \qquad 1\leq k<p.                                      \tag{11.4}
\end{equation}
Put $C_n=A_n\oplus B_n$.  Then $C_n\to A$ in
$\operatorname{Hom}(V,\mathcal S_p)$, while (11.1) and (11.4) give exact identity
\[
 D_{C_n}=e^P D_{p,C_n}\longrightarrow e^P D_{p,A}.
\]
This proves reverse inclusion in (DC), homeomorphism assertion, and principal fibre.
Pullback, unitary invariance, and direct-sum naturality follow directly from construction and
$D_{p,A\oplus B}=D_{p,A}D_{p,B}$.

\subsection*{Families and root-free cyclic blocks}

For Theorem~\ref{thm:parametrized-completion} roots must vary without branch choices.  Replace
root-of-unity block by weighted cyclic shift.  Given $k\geq1$, coefficient $c$, and linear form
$\ell$, let $C_{k,c,N}(v)$ act on $\C^k$ by
\[
 e_r\longmapsto a_r\ell(v)e_{r+1\ {\rm mod}\ k},\qquad
 a_0=\cdots=a_{k-2}=\delta,\quad
 a_{k-1}=\frac{c}{kN\delta^{k-1}},\quad
 \delta=N^{-1/k}.
\]
First string is empty when $k=1$.
For $0<q<k$, cyclic shift by $q$ has no diagonal entries, while
\[
 C_{k,c,N}(v)^k=\frac{c}{kN}\ell(v)^kI_k.
\]
Therefore $N$ copies have all lower trace moments zero and $k$th trace exactly
$c\ell(v)^k$.  No root of $c$ is chosen, and matrix entries depend linearly on $c$.  Singular
values of weighted cyclic shift are $|a_r\ell(v)|$.  On bounded coefficient set,
\begin{equation}
 \norm{\underbrace{C_{k,c,N}\oplus\cdots\oplus C_{k,c,N}}_{N\ {\rm copies}}}_p^p
 \leq C_{k,\ell,c}\,N^{1-p/k}\longrightarrow0.            \tag{11.5}
\end{equation}

Let now $c=c(x)$ vary continuously on compact $X$.  It is bounded, cyclic block varies
continuously, and (11.5) is uniform in $x$.  Choose once for each $k$ pure-power basis
$\{\ell_{\nu,k}^k\}_\nu$ of $\operatorname{Sym}^k(V^*)$.  Coefficients of every continuous
$P_k:X\to\operatorname{Sym}^k(V^*)$ in this basis are continuous.  Repeat ascending-degree
recursion uniformly on $X$.  Finite compressions $E_nA_xE_n$ converge uniformly because image
of compact $X$ is compact subset of $\operatorname{Hom}(V,\mathcal S_p)$.  This proves
sufficiency in (1.P).  Necessity and homeomorphism follow from normalized logarithm uniformly
near zero and continuous recovery of finitely many Taylor coefficients.

If $X$ is complex manifold and coefficients are holomorphic, weighted cyclic blocks are
holomorphic because entries are linear in $c(x)$.  Exhaust $X$ by compact sets and choose
multiplicities at stage $n$ large enough on first $n$ compacts.  Same diagonal argument gives
compact-open holomorphic statement and compatibility with restrictions and pullbacks.

\subsection*{Sharp capacity and statistics}

Let $B$ have rank $R$.  Schatten H\"older inequality and comparison of finite-dimensional
$\ell^k$ and $\ell^p$ norms give
\[
 |\Tr(B^k)|\leq\norm{B^k}_1
 \leq\norm{B}_k^k
 \leq R^{1-k/p}\norm{B}_p^k,
\]
proving (1.E).  If $R=kN$, diagonal operator with $N$ copies of
$t,t\omega,\ldots,t\omega^{k-1}$ has lower moments zero,
$\Tr(B^k)=Rt^k$, and equality throughout.  Solving equality for $R$ proves (1.R), with harmless
rounding to multiple of $k$.  At threshold $k=p$,
$|\Tr(B^p)|\leq\norm{B}_p^p$, so no nonzero moment survives $B\to0$ in $\mathcal S_p$.

It remains to prove exact statistics table.  Complex row follows from degree-by-degree blocks
above with $V=\C$.  If $B_n\geq0$, convergence of first moment makes $\Tr(B_n)$ bounded and
$\norm{B_n}_\infty\to0$; hence
\[
 0\leq\Tr(B_n^k)\leq\norm{B_n}_\infty^{k-1}\Tr(B_n)\longrightarrow0
 \quad(k\geq2).
\]
Every $m_1\geq0$ is realized by $N$ copies of $m_1/N$.

If $B_n=B_n^*$ and $p\geq3$, convergence of second moment makes $\Tr(B_n^2)$ bounded.  For
$k\geq3$,
\[
 |\Tr(B_n^k)|\leq\Tr|B_n|^k
 \leq\norm{B_n}_\infty^{k-2}\Tr(B_n^2)\longrightarrow0.
\]
Here $m_2\geq0$.  Any $m_1\in\R$ is realized by $N$ copies of $m_1/N$.  Any $m_2\geq0$ is
realized, without first moment, by $N$ copies each of
$\pm\sqrt{m_2/(2N)}$; its $\mathcal S_p$ norm tends to zero for $p>2$.  Direct sums give pair
$(m_1,m_2)$.  For $p=2$ only first coordinate exists.

Finally take positive blocks in even and odd degree and use supertrace.  For target coordinate
$c$ of degree $k$, put
\[
 w_{k,r}=(-1)^{k-r}\binom kr,\qquad 0\leq r\leq k.
\]
Use $N|w_{k,r}|$ copies of positive eigenvalue $t(r+1)$, placed in parity sign of $w_{k,r}$.
Then
\[
 \operatorname{Str}(B^q)=Nt^q
 \sum_{r=0}^k(-1)^{k-r}\binom kr(r+1)^q
 =
 \begin{cases}0,&q<k,\\Nt^kk!,&q=k.\end{cases}
\]
Choose parity reversal for sign of $c$ and
$t=(|c|/(Nk!))^{1/k}$.  Total $\mathcal S_p^p$ cost is
\[
 Nt^p\sum_{r=0}^k|w_{k,r}|(r+1)^p=O(N^{1-p/k})\longrightarrow0.
\]
Ascending-degree recursion realizes arbitrary real jet.  This proves
Theorem~\ref{thm:micro-spectrum} and completion law.

\begin{proof}[Proof of Corollary~\ref{cor:summability-phase}]
For real $r\geq1$, integers below $r$ are exactly $1,\ldots,m-1$ where $m=\lceil r\rceil$.
Finite-rank identity (11.1) holds with $m$ in place of $p$, and $\det_m$ is continuous under
$\mathcal S_r$ convergence because $\mathcal S_r\hookrightarrow\mathcal S_m$ continuously.
Block cost (11.3) becomes
\[
 (kN)^{1-r/k}|c|^{r/k}|\ell_i|^r,
\]
which tends to zero exactly for $k<r$.  Necessity proof is unchanged.  Rank--energy inequality
uses finite-dimensional norm comparison
$\norm B_k\leq R^{1/k-1/r}\norm B_r$ and gives exponent $r/(r-k)$.

For integer $p<q$ and $A\in\mathcal S_p$, all $A(v)^j$ with $j\geq p$ are trace class and
$T_{p,q}(A)$ is continuous polynomial in $A$.  Definitions give
\[
 D_{q,A}=D_{p,A}\exp T_{p,q}(A).
\]
This proves (1.S).  Identity
$T_{p,r}=T_{p,q}+T_{q,r}$ proves composition, and new quotient of jet spaces is
$J_{<p+1}(V)/J_{<p}(V)=\operatorname{Sym}^p(V^*)$.
\end{proof}

\section{Novelty boundary and significance}
\label{sec:boundary}

Following ingredients are established and are not claimed as contributions:
\begin{itemize}
\setlength{\itemsep}{0.35em}
\item definition, finite-rank trace formula, Schatten continuity, direct-sum property, and
      transition identity between $\det_p$ and $\det_{p+1}$ for regularized Fredholm
      determinants~\cite{Simon2005};
\item formula (1.1), finite-set density, extremal reduction to initial prime segments, and RH
      equivalence for generalized Euler constants~\cite{DiamondFord2008};
\item finite-set generalized Euler--Kronecker constants over number fields
      \cite{KandhilLunia2024}, and number-field Mertens asymptotics
      \cite{Rosen1999,GarciaLee2022};
\item fixed-operator factorization of normalized zeta determinant into exponential corrected
      residue polynomial and regularized Fredholm determinant~\cite{HartmannLesch2022};
\item Quillen and Bismut--Freed determinant lines, metrics, and connections
      \cite{Quillen1985,BismutFreed1986};
\item Quillen's polynomial-gauge picture and Dang's classification of local elliptic
      renormalized determinants as torsor under local polynomial counterterms
      \cite{Dang2022};
\item Wodzicki corrections to regularized Chern forms and determinant-bundle curvature
      \cite{PaychaRosenberg2000};
\item canonical traces, multiplicative anomalies, and residue determinants
      \cite{KontsevichVishik1994,Scott2005};
\item relative zeta determinants, analytic torsion, and determinant-line refinements
      \cite{Scott1999,Muller1998,DaiFreed1994,BravermanKappeler2007}.
\end{itemize}

New parent statement is Theorem~\ref{thm:completion-law}.  It does not claim invention of
$\det_p$ or identity (11.1).  New content is exact closure of \emph{ordinary finite-rank
determinant graph} under Schatten completion, coordinate-free multivariable fibre
$\exp(\bigoplus_{k<p}\operatorname{Sym}^kV^*)$, homeomorphic principal-extension structure, and
surjectivity by vanishing micro-spectra.  Hadamard factorization says how one entire function
with prescribed divisor differs from another; it does not prove which exponential factors arise
as simultaneous limits of finite-dimensional determinants while operator pencil converges in
$\mathcal S_p$.  Dang's theorem instead fixes elliptic background and classifies determinant-like
functions obtainable by \emph{local counterterm subtraction}; it neither takes graph closure in
Schatten topology nor classifies raw finite-rank approximants.  Conversely, our theorem imposes
no microlocal locality and does not replace Dang's local classification.  It proves that the
unconstrained Schatten completion already has precisely Quillen's principal polynomial-gauge
shape, with degree cutoff and every fibre direction forced by topology.
Theorem~\ref{thm:micro-spectrum} adds optimal rank--energy exponent and exact
positive/self-adjoint/graded-positive classification.  In particular, it explains why graded
analytic torsion is required to expose higher real defect directions.

New general statement is Theorem~\ref{thm:universal-law}: under $\mathcal S_p$ convergence,
corrected finite-part residue jet is complete cofinal discontinuity invariant.  It identifies
compact-open determinant cluster set homeomorphically, forces converse convergence of every jet
coefficient, fixes common divisor, and classifies divisor-cancelled torsion by polynomial
characters.  This is not fixed-operator factorization: it is stability-and-defect theorem for
sequences whose zeta determinants need not converge although operators do.

Theorem~\ref{thm:full-jet} proves sharpness in every Schatten order.  Its graded spectral
suspension is not formal repetition of scalar example: finite-difference weights annihilate all
lower zeta coordinates exactly, diagonal corrected finite parts remain independent, and a Baire
argument produces full-dimensional jet spectrum while $\mathcal S_p$ limits, integer residues,
and virtual divisor are fixed.  Hence universal law has exactly $p-1$ independent secondary
directions and full exponential-polynomial torsion group.

New arithmetic statements are Theorems~\ref{thm:classification} and
\ref{thm:numberfield-main}: arbitrary finite-set density does not imply, and does not contain,
classification under nested exhaustive ordering.  Prime-ideal relative extremality creates a
new universal lower endpoint $e^{-\gamma}$ over every number field.  Cancellation
$\rho_K\cdot e^{-\gamma}/\rho_K=e^{-\gamma}$ is forced only after finite part and spectral
residue are combined.  Alternating natural-fill and remote prime-ideal-block controls realize
every interval above endpoint.  Determinant lift and exact prime-ideal excision turn this
order-theoretic classification into sharp analytic no-descent theorem.

Distinction is not cosmetic.  Regularized Fredholm determinant converges in Hilbert--Schmidt
topology, primary residue converges, and divisor converges, but relative zeta determinant need
not converge.  Its entire failure is measured by one finite-part coordinate, and
Theorem~\ref{thm:classification} determines every possible cluster geometry of that coordinate.
Plain determinant-line triviality is insufficient in strongest possible sense for this model: it
forgets a complete real relative-torsion group.

Theorem~\ref{thm:generic} strengthens existence classification: maximal torsion ambiguity is
Baire-generic, every smaller ambiguity remains dense, and no ordering is point of continuity for
cluster spectrum.  Thus instability cannot be removed by selecting a topologically typical
cofinal ordering.

\section{Conclusion}

Parent universal law is determinant completion:
\[
 \overline{\operatorname{Graph}(\det_{\rm fin})}^{\,
   \operatorname{Hom}(V,\mathcal S_p)\times\operatorname{Hol}_1(V)}
 \cong
 \operatorname{Hom}(V,\mathcal S_p)\times
 \exp\!\left(\bigoplus_{k=1}^{p-1}\operatorname{Sym}^k(V^*)\right).
\]
Schatten topology preserves regularized core and forgets exactly low trace jet.  Arbitrarily
small finite-rank micro-spectra realize every forgotten coordinate, with sharp rank cost
$\varepsilon^{-pk/(p-k)}$ in degree $k$.  Positive spectra collapse hierarchy to degree one;
self-adjoint spectra to degrees one and two; superdeterminant group completion restores all
real degrees.  Theorem~\ref{thm:parametrized-completion} shows fibre persists functorially over
continuous and holomorphic bases, without root choices.  Corollary~\ref{cor:summability-phase}
identifies integer summability walls and associated graded anomaly layers.  These laws are
independent of zeta functions.

Zeta-determinant shadow is
\[
 L_n^{-1}\to A\text{ in }\mathcal S_p
 \quad\Longrightarrow\quad
 \operatorname{Clust}_{\rm co}(\mathcal D_n)
 =\left\{e^{P_b}\det{}_p(I-zA):
 b\in\operatorname{Clust}(b(L_n))\right\}.
\]
Limiting operator fixes divisor-carrying core.  Corrected finite-part residue jet fixes entire
remaining ambiguity, and relative determinant torsion is exponential polynomial of jet
difference.  Wodzicki residues supply local part of jet; finite parts supply nonlocal part.
Neither determinant-line triviality nor residue data alone is complete.

Finite-difference spectral suspensions show bound is sharp: in Schatten order $p$, exactly
$p-1$ independent corrected coordinates can survive, and a Baire-generic graded arithmetic
family realizes cluster set with nonempty interior in $\R^{p-1}$.  Relative torsions then generate
all real exponential-polynomial characters of degree below $p$.

Prime-ideal arithmetic realizes law sharply.  Dropping Euler factors creates two analytically
different boundary data at $s=1$:
\[
 \Res_{s=1}Z_{K,P}(s)=r_K(P),\qquad
 \FP_{s=1}Z_{K,P}(s)=\kappa_K(P).
\]
First supplies logarithmic cutoff correction; second survives as determinant torsion.  Along
every prime-ideal ordering, first vanishes and Hilbert--Schmidt operator tends to rank one.
Second has exactly closed-interval cluster spectra in $[e^{-\gamma},\infty]$, with endpoint
independent of number field.  Limiting divisor is rigid but zeta determinant is a real torsor:
\[
 \Delta_c(z)=e^{(1-c)z}(1-z),\qquad
 \frac{\Delta_c}{\Delta_d}=e^{(d-c)z}.
\]
Thus natural arithmetic cofinal systems realize every interval of complete defect coordinate,
while maximal interval and full real torsion group are Baire-generic.

\begin{thebibliography}{99}

\bibitem{BismutFreed1986}
J.-M.~Bismut and D.~S.~Freed,
\emph{The analysis of elliptic families. II. Dirac operators, eta invariants, and the holonomy
theorem},
Communications in Mathematical Physics \textbf{107} (1986), 103--163.

\bibitem{BravermanKappeler2007}
M.~Braverman and T.~Kappeler,
\emph{Refined analytic torsion as an element of the determinant line},
Geometry \& Topology \textbf{11} (2007), 139--213,
\href{https://doi.org/10.2140/gt.2007.11.139}{doi:10.2140/gt.2007.11.139}.

\bibitem{Bunke2002}
U.~Bunke,
\emph{Index theory, eta forms, and Deligne cohomology},
arXiv:math/0201112,
\url{https://arxiv.org/abs/math/0201112}.

\bibitem{DaiFreed1994}
X.~Dai and D.~S.~Freed,
\emph{Eta-invariants and determinant lines},
Journal of Mathematical Physics \textbf{35} (1994), 5155--5194,
\url{https://arxiv.org/abs/hep-th/9405012}.

\bibitem{Dang2022}
N.~V.~Dang,
\emph{Renormalization of determinant lines in quantum field theory},
Analysis \& PDE \textbf{15} (2022), 1--62,
\href{https://doi.org/10.2140/apde.2022.15.1}{doi:10.2140/apde.2022.15.1}.

\bibitem{DiamondFord2008}
H.~G.~Diamond and K.~Ford,
\emph{Generalized Euler constants},
Mathematical Proceedings of the Cambridge Philosophical Society \textbf{145} (2008), 27--41,
\href{https://doi.org/10.1017/S0305004108001187}{doi:10.1017/S0305004108001187}.

\bibitem{GarciaLee2022}
S.~R.~Garcia and E.~S.~Lee,
\emph{Unconditional explicit Mertens' theorems for number fields and Dedekind zeta residue bounds},
Ramanujan Journal \textbf{57} (2022), 1169--1191,
\href{https://doi.org/10.1007/s11139-021-00435-6}{doi:10.1007/s11139-021-00435-6}.

\bibitem{HartmannLesch2022}
L.~Hartmann and M.~Lesch,
\emph{Zeta and Fredholm determinants of self-adjoint operators},
Journal of Functional Analysis \textbf{283} (2022), 109491,
\href{https://doi.org/10.1016/j.jfa.2022.109491}{doi:10.1016/j.jfa.2022.109491}.

\bibitem{KandhilLunia2024}
N.~Kandhil and R.~Lunia,
\emph{Transcendence of generalised Euler--Kronecker constants},
Bulletin of the Australian Mathematical Society \textbf{109} (2024), 464--470,
\href{https://doi.org/10.1017/S0004972723000631}{doi:10.1017/S0004972723000631}.

\bibitem{KontsevichVishik1994}
M.~Kontsevich and S.~Vishik,
\emph{Determinants of elliptic pseudo-differential operators},
arXiv:hep-th/9404046,
\url{https://arxiv.org/abs/hep-th/9404046}.

\bibitem{Muller1998}
W.~Mueller,
\emph{Relative zeta functions, relative determinants and scattering theory},
Communications in Mathematical Physics \textbf{192} (1998), 309--347,
\href{https://doi.org/10.1007/s002200050301}{doi:10.1007/s002200050301}.

\bibitem{PaychaRosenberg2000}
S.~Paycha and S.~Rosenberg,
\emph{Curvature on determinant bundles and first Chern forms},
arXiv:math/0009172,
\url{https://arxiv.org/abs/math/0009172}.

\bibitem{PaychaScott2007}
S.~Paycha and S.~Scott,
\emph{A Laurent expansion for regularized integrals of holomorphic symbols},
Geometric and Functional Analysis \textbf{17} (2007), 491--536,
\href{https://doi.org/10.1007/s00039-007-0597-8}{doi:10.1007/s00039-007-0597-8}.

\bibitem{Quillen1985}
D.~Quillen,
\emph{Determinants of Cauchy--Riemann operators over a Riemann surface},
Functional Analysis and Its Applications \textbf{19} (1985), 31--34.

\bibitem{RaySinger1971}
D.~B.~Ray and I.~M.~Singer,
\emph{$R$-torsion and the Laplacian on Riemannian manifolds},
Advances in Mathematics \textbf{7} (1971), 145--210,
\href{https://doi.org/10.1016/0001-8708(71)90045-4}{doi:10.1016/0001-8708(71)90045-4}.

\bibitem{Rosen1999}
M.~Rosen,
\emph{A generalization of Mertens' theorem},
Journal of the Ramanujan Mathematical Society \textbf{14} (1999), 1--19.

\bibitem{Scott1999}
S.~Scott,
\emph{Relative zeta determinants and the Quillen metric},
arXiv:math/9910148,
\url{https://arxiv.org/abs/math/9910148}.

\bibitem{Scott2005}
S.~Scott,
\emph{The residue determinant},
Communications in Partial Differential Equations \textbf{30} (2005), 483--507,
\url{https://arxiv.org/abs/math/0406268}.

\bibitem{Simon2005}
B.~Simon,
\emph{Trace Ideals and Their Applications}, second edition,
Mathematical Surveys and Monographs \textbf{120}, American Mathematical Society, 2005.

\end{thebibliography}

\end{document}
